% !TeX root = ../../Book2.tex
% 纲要标题：## 附录 D　算子代数：对合、拓扑与表示
% 纲要位置：工作记录/计划纲要.md，第 1860 行。

% 纲要要点（供目录核对）：
% > 本附录联系代数与泛函分析，区分对合、范数完备化与算子拓扑闭包。在已列分析学输入下证明 Gelfand 对偶、连续函数演算、GNS、Gelfand--Naimark 表示及含幺自伴版本的双交换子定理；谱半径与解析谱映射、一般张量范数极值及自由群 II_1 因子另明外部归属。简要介绍投影类型，不展开完整分类。

\chapter{算子代数：对合、拓扑与表示}
\label{b2:app:D}

\begin{chapterabstract}
在复代数上加入对合、范数与完备性以后，模表示的问题转为 Hilbert 空间上的有界表示问题。本附录由矩阵、连续函数和紧算子引入这些结构，说明交换 Gelfand 对偶的含义，构造态的 GNS 表示，并证明含幺自伴算子代数的双交换子定理。投影的等价与有限性给出因子的 I、II、III 型；整数群和矩阵张量积则提供可计算的比较对象，帮助区分代数有限和、范数完备化与弱算子闭包。
\end{chapterabstract}

\begin{learninggoals}
\begin{itemize}
\item 核对对合、范数、完备性及单位元各自的作用，区分有界表示与 Weyl 代数的代数作用。
\item 理解 Gelfand 变换证明中的谱、紧性与逼近步骤，并在有限集合和正规矩阵上计算函数演算。
\item 从正泛函构造内积商空间，证明左作用有界并完成 GNS 表示。
\item 区分三种算子收敛，掌握有限次放大在双交换子证明中的用途，并用投影判据区分 I、II、III 型因子。
\item 比较全群、约化群与群 von Neumann 代数，以及代数张量积与其完备化。
\end{itemize}
\end{learninggoals}

在 \([0,1]\) 上，多项式可以逐项相加、相乘，也能按一致范数相互逼近。一列多项式的一致极限却可能是 \(|t-1/2|\) 这样的非多项式函数。若关心极限后仍可进行的运算，就必须把函数空间完备化，并让复共轭与范数一同进入结构。所得对象不再只由有限个代数等式描述，连续性成为判断表示是否合适的一部分。

\ProseParagraphBreak

矩阵代数和连续函数代数分别提供非交换与交换的基本模型。前者让我们计算伴随、正性及向量态，后者把点的取值同代数上的乘法泛函联系起来。一般抽象代数能否在 Hilbert 空间上这样实现，要经过态、内积和完备化的构造。实现为有界算子以后，还须选择哪些极限留在代数中：范数收敛控制所有单位向量，强算子收敛逐个控制向量，弱算子收敛只控制矩阵系数。对含幺自伴子代数，后两种闭包将由双交换子定理识别为同一个双交换子，范数闭包却可能更小。整数群的 Laurent 多项式最后会分别完成为圆周连续函数和本质有界函数，给出这一差别的具体模型。

\ProseParagraphBreak

本附录使用 Banach 与 Hilbert 空间、有界算子及伴随、闭子空间的正交投影、完备化和紧 Hausdorff 空间的基础。谱与态的论证还要用 Liouville 定理\footnote{刘维尔（Joseph Liouville，1809-1882），法国数学家，研究分析、数论与微分方程。}、Hahn--Banach 保范延拓定理\footnote{哈恩（Hans Hahn，1879-1934），奥地利数学家，研究泛函分析与集合论。}及紧空间的乘积定理。谱半径公式与解析谱映射在\secref{b2:sec:D02}中准确陈述，证明所需的复变积分参见所引谱论文献。函数逼近使用 Stone--Weierstrass 定理\footnote{斯通（Marshall Harvey Stone，1903-1989），美国数学家，研究算子理论与函数逼近。魏尔斯特拉斯（Karl Theodor Wilhelm Weierstrass，1815-1897），德国数学家，奠定严格的实分析基础，并证明多项式逼近定理。}；圆周例子另用 Fourier--Plancherel 定理\footnote{傅里叶（Joseph Fourier，1768-1830），法国数学家及物理学家，建立傅里叶级数理论，研究热传导。普朗歇雷尔（Michel Plancherel，1885-1967），瑞士数学家，证明傅里叶变换的等距定理。}。内积沿用本卷 Hermitian 型的约定：第一变量共轭线性、第二变量线性。

\ProseParagraphBreak
从这些分析结果出发，连续函数演算使谱上的非负函数变成代数中的正元素；态由此给出内积，GNS 构造再将左乘变成有界算子。另一组问题从算子收敛出发，研究交换关系如何确定弱闭包。两组论证的主要依赖见\cref{b2:fig:operator-dependencies}。在已列出的分析学输入下，本附录展开 Gelfand 对偶、连续函数演算、GNS、Gelfand--Naimark 表示及当前含幺版本的双交换子证明；谱论基础引理只证明其中的基本部分，谱半径公式与解析谱映射按所引文献使用。自由群 \(L(F_2)\) 的 \(\mathrm{II}_1\) 型结构和一般张量范数极值则引用外部理论。张量范数的极值定理还需纯态与不可约表示、态空间的弱星紧性、紧凸集的极点与分离，以及非含幺代数的近似单位；\secref{b2:sec:D05}引用该定理，并说明含幺情形的最大性。

\begin{figure}[htbp]
\centering
% !TeX root = ../Book2.tex
\begin{tikzpicture}[
  concept/.style={draw, rounded corners=2pt, align=center,
    text width=20mm, minimum height=11mm, inner sep=2pt, font=\small},
  every path/.style={semithick}]
\node[concept] (cstar) at (0,0) {\(C^*\)-代数\\与范数};
\node[concept] (spectrum) at (2.6,0) {谱论与\\Gelfand 对偶};
\node[concept] (calculus) at (5.2,0) {连续\\函数演算};
\node[concept] (positive) at (7.8,0) {正性\\与态};
\node[concept] (gns) at (10.4,0) {GNS\\构造};
\node[concept] (faithful) at (13,0) {忠实\\Hilbert 表示};
\draw[->] (cstar) -- (spectrum);
\draw[->] (spectrum) -- (calculus);
\draw[->] (calculus) -- (positive);
\draw[->] (positive) -- (gns);
\draw[->] (gns) -- (faithful);
\node[concept] (operators) at (0,-1.7) {\(B(H)\)};
\node[concept] (topology) at (2.6,-1.7) {SOT 与 WOT};
\node[concept] (bicommutant) at (5.2,-1.7) {双交换子\\定理};
\node[concept] (group) at (7.8,-1.7) {群 von Neumann\\代数};
\node[concept] (extra) at (10.4,-1.7) {纯态、极点\\与近似单位};
\node[concept] (tensor) at (13,-1.7) {张量范数\\极值定理};
\draw[->] (operators) -- (topology);
\draw[->] (topology) -- (bicommutant);
\draw[->] (bicommutant) -- (group);
\draw[->] (faithful) -- node[right,font=\footnotesize] {空间模型} (tensor);
\draw[->] (extra) -- (tensor);
\end{tikzpicture}

\caption[算子代数的主要依赖]{算子代数的主要依赖。箭头表示论证所用的结构或结果；张量范数极值还需要图中另列的分析学输入。}
\label{b2:fig:operator-dependencies}
\end{figure}

进一步阅读可用 Blackadar 的专著《Operator Algebras》\cite{Blackadar2006OperatorAlgebras}。本附录的具体编号依据作者的 2017 年校订本\cite[\nopp I.9.1; II.2; II.6.4; II.9; III.1]{Blackadar2017OperatorAlgebrasCorrected}。

% !TeX root = ../../Book2.tex
% 纲要标题：### D.1 从带对合代数到 \(C^*\)-代数
% 纲要位置：工作记录/计划纲要.md，第 1673 行。

% 纲要要点（供目录核对）：
% * 复代数上的共轭线性对合、Banach 代数及 \(C^*\) 恒等式 \(\|a^*a\|=\|a\|^2\)
% * 例子：\(M_n(\mathbb C)\)、\(C_0(X)\)、\(B(H)\) 与紧算子代数；说明单位元和完备性的不同情况
% * 抽象环、带对合代数和算子代数的附加结构不能省略；Weyl/Ore 微分算子环不因此自动成为 \(C^*\)-代数

\section{从带对合代数到 \texorpdfstring{\(C^*\)}{C*}-代数}
\label{b2:sec:D01}

矩阵既可以相乘，也可以取共轭转置、作用于带内积的空间并计算算子范数。这些附加结构彼此有关，却不能从抽象环的乘法中一并省略。以下讨论复代数；Hilbert 空间的内积沿用\secref{b2:sec:1004}的约定，对第一变量共轭线性、第二变量线性。所有表示均由有界算子给出，除非明确另作说明。

\subsection{对合、Banach 代数与 \texorpdfstring{\(C^*\)}{C*} 恒等式}
\label{b2:subsec:D01:01}

在带内积的空间中，\(\langle T\xi,T\xi\rangle=\langle\xi,T^*T\xi\rangle\) 将向量长度与乘积 \(T^*T\) 联系起来。算子范数由所有单位向量的长度变化确定，因此对合与范数不能任意搭配。为了在抽象代数中保留这一关系，还要让范数极限留在代数内，并规定下面的 \(C^*\) 恒等式。

\begin{definition}\label{b2:def:banach-star-cstar}
设 \(A\) 为复结合代数。若映射 \(a\mapsto a^*\) 对任意 \(a,b\in A\)、\(\lambda\in\mathbb C\) 满足
\[
(a+b)^*=a^*+b^*,\quad (\lambda a)^*=\bar\lambda a^*,\quad
(ab)^*=b^*a^*,\quad (a^*)^*=a.
\]
则称它为 \(A\) 上的\term[gonge-xianxing-duihe]{共轭线性对合}[conjugate-linear involution]。若 \(A\) 带有使其完备、并满足 \(\|ab\|\le\|a\|\cdot\|b\|\) 的范数，则称 \(A\) 为\term[Banach-daishu]{Banach 代数}[Banach algebra]。若 \(A\) 还是带有上述对合的 Banach 代数，并对任意 \(a\in A\) 满足
\begin{equation}\label{b2:eq:cstar-identity}
\|a^*a\|=\|a\|^2\qquad(a\in A),
\end{equation}
则称 \(A\) 为 \term[Cxing-daishu]{\(C^*\)-代数}[C*-algebra]。这里不默认存在单位元。若两个带对合的复代数之间的映射保持乘法、复线性结构和对合，则称它为 \term[xing-tongtai]{\(*\)-同态}[*-homomorphism]；需要保持单位元时另行注明。
\end{definition}

由 \(\|a\|^2=\|a^*a\|\le\|a^*\|\cdot\|a\|\)，对 \(a\ne0\) 消去范数得到 \(\|a\|\le\|a^*\|\)；再把 \(a\) 换为 \(a^*\)，可得 \(\|a^*\|=\|a\|\)。因此对合自动连续。非零含幺情形还给出 \(1^*=1\)、\(\|1\|=1\)。这里共轭线性的含义尤其要与附录B中固定基域的线性对合区别：复数标量 \(i\) 被送到 \(-i\)。

\subsection{矩阵、函数与算子代数}
\label{b2:subsec:D01:02}

设 \(H\ne0\) 是复 Hilbert 空间。全部有界线性算子组成 \(B(H)\)，乘法为复合，对合为伴随，范数为
\[
\|T\|=\sup_{\|\xi\|\le1}\|T\xi\|.
\]
\symidx{\(B(H)\)：Hilbert 空间上的有界线性算子代数}
有界算子空间的完备性是这里所需的分析学基础。伴随满足 \(\langle T\xi,\eta\rangle=\langle\xi,T^*\eta\rangle\)，并由内积的\term[Cauchy-Schwarz-budengshi]{Cauchy--Schwarz 不等式}[Cauchy--Schwarz inequality]及单位向量上的上确界得到 \(\|T^*\|=\|T\|\)。

\begin{proposition}\label{b2:prop:operator-cstar}
设 \(H\ne0\) 为复 Hilbert 空间，\(B(H)\) 为其有界复线性算子代数，对合取 Hilbert 空间伴随，则 \(B(H)\) 在算子范数下是含幺 \(C^*\)-代数；它的每个范数闭且对伴随封闭的复子代数也是 \(C^*\)-代数，未必含有 \(I_H\)。
\end{proposition}
\begin{proof}
次乘性来自 \(\|ST\xi\|\le\|S\|\cdot\|T\xi\|\)。对单位向量 \(\xi\)，
\[
\|T\xi\|^2=\langle\xi,T^*T\xi\rangle\le\|T^*T\|,
\]
故 \(\|T\|^2\le\|T^*T\|\)。反向由次乘性及 \(\|T^*\|=\|T\|\) 得到。闭子代数继承完备性、对合和范数恒等式，所以也满足定义。
\end{proof}

取 \(H=\mathbb C^n\)，即得 \(M_n(\mathbb C)\)；其范数是 Euclid 内积诱导的算子范数，不能任意改用另一种矩阵范数。例如把算子范数统一乘二，所得范数 \(N\) 仍完备且满足次乘性，但对 \(E_{11}\) 有
\[
N(E_{11}^*E_{11})=N(E_{11})=2,
\qquad N(E_{11})^2=4.
\]
所以它不满足\eqref{b2:eq:cstar-identity}。这个恒等式确定了对合与范数之间的数值关系，单有完备性和次乘性还不够。

\ProseParagraphBreak
设 \(X\) 为局部紧 Hausdorff 空间。记 \(C_0(X)\) 为连续复函数 \(f\) 的代数，要求对每个 \(\varepsilon>0\)，集合 \(\Bigl\{\,x\in X\,\Bigm|\,\bigl|f(x)\bigr|\ge\varepsilon\,\Bigr\}\) 紧。这是“在无穷远处趋于零”的精确定义。对合为逐点复共轭，范数为上确界范数；一致极限保持上述性质，故它完备，而 \(\|\bar f f\|_\infty=\|f\|_\infty^2\)。当 \(X\) 紧时，它就是含有常值函数一的 \(C(X)\)；\(C_0(\mathbb R)\) 则没有单位元。若它有单位 \(e\)，对每个点选一个在该点非零的紧支集连续函数，\(ef=f\) 迫使 \(e\) 处处为一，与趋零条件矛盾。
\symidx{\(C_0(X)\)：在无穷远处趋于零的连续复函数代数}
\symidx{\(C(X)\)：紧 Hausdorff 空间上的连续复函数代数}

\ProseParagraphBreak
有限秩算子组成 \(B(H)\) 的双边 \(*\)-理想；其范数闭包记为 \(\mathcal K(H)\)，称为\term[jinsuanzi-daishu]{紧算子代数}[algebra of compact operators]。在 Hilbert 空间上，这与把有界集送到相对紧集的算子一致。确实，有限秩算子的单位球像在有限维空间中全有界；范数极限仍使单位球像全有界。反过来，若 \(T\) 的单位球像相对紧，取其有限 \(\varepsilon\)-网，让 \(P\) 为网中向量张成空间的正交投影，便有 \(\|T-PT\|\le\varepsilon\)，而 \(PT\) 有限秩。完备空间中全有界集的闭包紧，故两种定义等价，亦见 Blackadar\cite[\nopp I.8.1.1; I.8.1.5]{Blackadar2006OperatorAlgebras}。

\ProseParagraphBreak
若 \(H\) 无限维，任何有限秩 \(F\) 都有一个单位核向量，故 \(\|I-F\|\ge1\)，所以 \(I\notin\mathcal K(H)\)。它也没有别的单位：若 \(e\) 对全部秩一算子作左单位，逐一作用于其像向量就得到 \(e\xi=\xi\) 对所有 \(\xi\)，从而 \(e=I\)。
\symidx{\(\mathcal K(H)\)：Hilbert 空间上的紧算子代数}

\subsection{抽象环、对合代数与算子代数}
\label{b2:subsec:D01:03}

上三角矩阵代数有天然算子范数，却一般不对共轭转置封闭；有限秩算子对伴随封闭，却在无限维时不完备。一个例子缺少对合封闭性，另一个缺少闭性，不能仅因都由算子组成就称它们为 \(C^*\)-代数。

\begin{proposition}\label{b2:prop:bounded-weyl-impossible}
设 \(A\ne0\) 为含幺复 Banach 代数，则不存在元素 \(x,d\in A\) 满足 \(dx-xd=1_A\)。因此，对任意非零复 Hilbert 空间 \(H\)，复 Weyl 代数 \(A_1(\mathbb C)\) 没有取值于 \(B(H)\) 的保幺代数表示。
\end{proposition}
\begin{proof}
若关系成立，归纳展开交换子得到 \(dx^n-x^nd=nx^{n-1}\)。于是
\[
n\|x^{n-1}\|\le 2\|d\|\cdot\|x^n\|\qquad(n\ge1).
\]
\(d\ne0\)，从 \(\|1\|>0\) 开始迭代，得到
\[
\|x\|^n\ge\|x^n\|\ge\frac{n!}{\bigl(2\|d\|\bigr)^n}\|1\|.
\]
但 \((n!)^{1/n}\) 无界，例如最后至少 \(n/2\) 个因子都不小于 \(n/2\)。取 \(n\) 次方根即得矛盾。
\end{proof}

第20章的求导与乘坐标作用并没有违反此命题：它们首先定义在多项式空间上，若置于通常的函数 Hilbert 空间，便涉及无界算子及共同定义域。处理这类作用需要另外的分析，而不能直接把 Ore 扩张的代数表示当作有界表示。

% !TeX root = ../../Book2.tex
% 纲要标题：### D.2 交换情形：谱与 Gelfand 对偶
% 纲要位置：工作记录/计划纲要.md，第 1679 行。

% 纲要要点（供目录核对）：
% * 特征标空间（非零复代数同态）和 Gelfand 变换；交换含幺 \(C^*\)-代数与紧 Hausdorff 空间的对偶
% * 非含幺情形的 \(C_0(X)\)；证明正规元连续函数演算的基本形式
% * 区分 Gelfand 谱及其拓扑与第三卷的素谱及 Zariski 拓扑，不把两种几何直接等同

\section{交换情形：谱与 Gelfand 对偶}
\label{b2:sec:D02}

对角矩阵可以视为有限集合上的函数。Gelfand 理论把这个观察推广到一般交换 \(C^*\)-代数：空间中的点由取值同态恢复，函数的一致收敛则对应代数范数。本节说明一般定理的证明方法，再把有限集合情形及映射方向算清楚。

\subsection{特征标空间与 Gelfand 变换}
\label{b2:subsec:D02:01}

\begin{definition}\label{b2:def:gelfand-spectrum}
设 \(A\ne0\) 为含幺交换复 Banach 代数。若映射 \(\chi:A\to\mathbb C\) 非零、复线性且保持乘法，则称 \(\chi\) 为 \(A\) 的\term[daishu-tezhengbiao]{代数特征标}[character of an algebra]。将全部特征标组成的集合 \(\Delta(A)\) 赋予使每个取值映射 \(\chi\mapsto\chi(a)\)、\(a\in A\)，连续的最弱拓扑，所得空间称为 \(A\) 的\term[Gelfand-pu]{Gelfand 谱}[Gelfand spectrum]。定义映射
\[
\Gamma:A\longrightarrow C\bigl(\Delta(A)\bigr),\qquad
\Gamma(a)(\chi)=\chi(a),
\]
并称它为 \(A\) 的\term[Gelfand-bianhuan]{Gelfand 变换}[Gelfand transform]。
\symidx{\(\Delta(A)\)：交换 Banach 代数的 Gelfand 谱}
\symidx{\(\Gamma\)：Gelfand 变换}
\end{definition}

非零性使 \(\chi(1)=1\)。若 \(|\lambda|>\|a\|\)，\term[Neumann-jishu]{Neumann 级数}[Neumann series] \(\sum_{j\ge0}(a/\lambda)^j\) 在 Banach 代数中收敛，并是 \(1-a/\lambda\) 的逆。\footnote{诺伊曼（Carl Gottfried Neumann，1832-1925），德国数学家，研究位势论与积分方程。此处与下文 von Neumann 代数所纪念的数学家并非同一人。}把可逆元送到可逆元可知 \(\chi(a)\ne\lambda\)，因而 \(\bigl|\chi(a)\bigr|\le\|a\|\)，故 \(\|\chi\|\le1\)。在 \(C^*\)-代数情形，\(\|1\|=1\)，进一步有 \(\|\chi\|=1\)。所以这些同态自动连续，Gelfand 变换的值确实是连续函数。这里的特征标是代数同态，与第五卷一般表示的迹特征标不是同一个定义。

\begin{theorem}[Gelfand 对偶]\label{b2:thm:gelfand-duality-external}
设 \(A\ne0\) 为含幺交换复 \(C^*\)-代数，\(\Delta(A)\) 为其代数特征标空间，\(\Gamma:A\rightarrow C(\Delta(A))\) 由 \(\Gamma(a)(\chi)=\chi(a)\) 定义，则 \(\Delta(A)\) 为紧 Hausdorff 空间，\(\Gamma\) 为等距 \(*\)-同构。反过来，对非空紧 Hausdorff 空间 \(X\)，\(C(X)\) 的全部特征标恰为取值映射 \(f\mapsto f(x)\)、\(x\in X\)，且 \(X\cong\Delta\bigl(C(X)\bigr)\) 为同胚。对应 \(X\leftrightarrow C(X)\) 将连续映射反向对应为保幺 \(*\)-同态。
\end{theorem}

\begin{proofsketch}
先说明所需的 Banach 代数谱论。对含幺 Banach 代数中的元素 \(a\)，不能处处定义逆 \((a-\lambda1)^{-1}\)：否则将任一有界线性泛函作用于这个逆，便得到一个在无穷远趋于零的整函数，由 Liouville 定理\footnote{刘维尔（Joseph Liouville，1809-1882），法国数学家，研究分析、数论与微分方程。}恒为零；有界线性泛函分离点，就与逆非零矛盾。因此不可逆的 \(\lambda\) 构成非空紧集。展开逆的 Neumann 级数，并比较其收敛半径，得到谱半径公式
\[
r(a)=\lim_{m\to\infty}\|a^m\|^{1/m}.
\]
若 \(h=h^*\)，则 \(e^{ith}\) 对每个实数 \(t\) 都酉，范数为一。指数函数的谱映射使 \(e^{it\lambda}\) 在 \(\lambda\) 属于 \(h\) 的谱时具有模一，故 \(\lambda\) 为实数。\(C^*\) 恒等式又给出 \(\|h^{2^m}\|=\|h\|^{2^m}\)，因此 \(r(h)=\|h\|\)；对正规元同样可由 \(a^*a=aa^*\) 得到 \(r(a)=\|a\|\)。这里省略复变积分对谱半径公式和指数谱映射的推导，见\cite[\nopp II.1.4--II.1.6]{Blackadar2006OperatorAlgebras}。

回到交换代数 \(A\)。每个极大理想 \(M\) 都闭：其闭包若含单位元，便含一个距单位元小于一的元素，而该元素由 Neumann 级数可逆，与理想为真矛盾。商 \(A/M\) 因而是 Banach 除代数；上述谱非空性说明每个商元素都等于某个复数，故 \(A/M\cong\mathbb C\)。于是极大理想给出特征标。更具体地，若 \(a-\lambda1\) 不可逆，它生成的真理想包含在某个极大理想内，所得特征标满足 \(\chi(a)=\lambda\)。所以谱恰为全部特征标值。

各特征标属于紧乘积 \(\prod_{a\in A}\bigl\{\,z\,\bigm|\,|z|\le\|a\|\,\bigr\}\)。线性、乘法和 \(\chi(1)=1\) 都是闭条件，故 \(\Delta(A)\) 是此乘积的闭子集，紧且 Hausdorff。自伴元的特征标值为实数，从而 \(\Gamma\) 保持对合；又由正规元的范数公式，
\[
\|\Gamma(a)\|_\infty=r(a)=\|a\|.
\]
其像因此闭，且含常数、分离点、对共轭封闭。复 Stone--Weierstrass 定理说明像稠密，于是像为全部连续函数。

最后，若 \(\chi\) 是 \(C(X)\) 的特征标，其核中函数必有共同零点。否则由紧性选有限个 \(f_j\in\ker\chi\)，使 \(\sum_j|f_j|^2\) 处处正而可逆，却仍属于核，矛盾。取共同零点 \(x\)，核中的 \(f-\chi(f)1\) 给出 \(\chi(f)=f(x)\)。连续函数分离点保证 \(x\) 唯一；紧空间到 Hausdorff 空间的连续双射给出所述同胚。对保幺 \(*\)-同态 \(\psi:C(Y)\to C(X)\)，\(\operatorname{ev}_x\psi\) 唯一确定一个点 \(f(x)\in Y\)，且 \(\psi(h)=h\circ f\)。取值拓扑使 \(f\) 连续，复合关系给出反变对应。紧乘积定理、Liouville 定理和 Stone--Weierstrass 逼近是这里所用的分析学基础；其在该论证中的用法及完整谱论细节见\cite[\nopp I.2; II.2.1--II.2.2, Theorems II.2.2.4, II.2.2.6]{Blackadar2006OperatorAlgebras}。
\end{proofsketch}

这称为\term[Gelfand-duiou]{Gelfand 对偶}[Gelfand duality]。“每个点给出同态”只是构造的一半，全部同态是否如此、Gelfand 变换的像是否穷尽，分别由上面的共同零点论证和逼近论证解决。

\ProseParagraphBreak
在 \(A=\mathbb C^r\) 中，这些步骤可以直接完成。记坐标幂等元为 \(e_i\)。\(\chi(e_i)^2=\chi(e_i)\) 使其值只能是零或一；正交性与 \(\sum_i\chi(e_i)=1\) 说明恰有一个为一。因此 \(\chi\) 就是取某个坐标，\(\Delta(A)\) 是 \(r\) 个点的离散空间。Gelfand 变换把坐标向量当作这些点上的函数，上确界范数恰为 \(\max_i|a_i|\)。

\ProseParagraphBreak
若 \(f:X\to Y\) 连续，则反向同态为 \(f^*:C(Y)\to C(X)\)、\(h\mapsto h\circ f\)。对 \(g:Y\to Z\)，复合满足 \((g\circ f)^*=f^*\circ g^*\)。这与第五章函子的反变约定一致；并不是把点沿同一方向线性延拓成一个函数。

\subsection{非含幺情形与连续函数演算}
\label{b2:subsec:D02:02}

若交换 \(C^*\)-代数 \(A\) 不含幺，仍取非零特征标，得到局部紧 Hausdorff 空间，相应同构是 \(A\cong C_0\bigl(\Delta(A)\bigr)\)。证明可归到含幺情形：在添幺代数 \(\widetilde A\) 上应用上一定理，商同态 \(\widetilde A\to\widetilde A/A\cong\mathbb C\) 对应一个点 \(\infty\)。其他特征标限制到 \(A\) 均非零，反过来非零特征标由 \(\widetilde\chi(a+\lambda1)=\chi(a)+\lambda\) 唯一延伸。故 \(\Delta(A)=\Delta(\widetilde A)\setminus\{\infty\}\)，而 \(A\) 对应在 \(\infty\) 为零的连续函数，恰为此局部紧空间上的 \(C_0\)。添幺代数的范数构造见\secref{b2:sec:D03}末的说明。

\ProseParagraphBreak
对任意连续 \(f:X\to Y\)，拉回未必把 \(C_0(Y)\) 送入 \(C_0(X)\)：例如从 \(\mathbb R\) 到一个点的常值映射，将常值一拉回为不趋零的函数。若 \(f\) 是真映射，即紧集的原像紧，拉回便保持趋零，因为 \(\{|h\circ f|\ge\varepsilon\}=f^{-1}(\{|h|\ge\varepsilon\})\) 紧。一般 \(*\)-同态 \(\psi:C_0(Y)\to C_0(X)\) 则在开集 \(U=\{x:\operatorname{ev}_x\psi\ne0\}\) 上给出取值点 \(f(x)\in Y\)，在 \(X\setminus U\) 上取值为零。开性来自某个 \(\psi(h)\) 在该点非零；添幺后的紧空间对应还说明 \(f:U\to Y\) 为真映射。反过来，真映射保证 \(h\circ f\) 的非零水平集紧包含于 \(U\)，故把它在外部补零就得到 \(C_0(X)\) 中的连续函数。这给出完整对应，见\cite[\nopp II.2.2.7--II.2.2.8]{Blackadar2006OperatorAlgebras}。

\ProseParagraphBreak
前面的矩阵计算已经允许把多项式 \(p\) 代入一个元素 \(a\)，得到 \(p(a)\)。函数代数说明还需要更广的运算：在 \(C([0,1])\) 中，令 \(h(t)=t\)，则 \(r(t)=\sqrt t\) 是连续函数，却不等于任何多项式 \(p(h)\)。否则多项式 \(p(t)^2-t\) 在整个区间上为零，由多项式根数界恒为零；比较次数却得到 \(2\deg p=1\)，矛盾。连续函数演算要使这种由连续函数产生的新元素也能由 \(a\) 确定。函数的定义域来自 \(a-\lambda1\) 不可逆的那些 \(\lambda\)，这便引出元素的谱。

\begin{definition}\label{b2:def:banach-element-spectrum}
设 \(A\ne0\) 为含幺复 Banach 代数，\(a\in A\)。定义集合
\[
\sigma_A(a)=\{\,\lambda\in\mathbb C:\ a-\lambda1\text{ 不可逆}\,\}.
\]
称它为 \(a\) 在 \(A\) 中的\term[yuansu-pu]{谱}[spectrum of an element]。
\symidx{\(\sigma_A(a)\)：元素在含幺 Banach 代数中的谱}
若 \(A\) 还是 \(C^*\)-代数且 \(a^*a=aa^*\)，则称 \(a\) 为\term[zhenggui-yuansu]{正规元}[normal element]。
\end{definition}

\begin{theorem}[连续函数演算]\label{b2:thm:continuous-functional-calculus}
设 \(A\ne0\) 为含幺复 \(C^*\)-代数，\(a\in A\) 满足 \(a^*a=aa^*\)，则其谱 \(\sigma_A(a)\) 为非空紧集，存在唯一的保幺等距 \(*\)-同构
\[
C\bigl(\sigma_A(a)\bigr)\longrightarrow C^*(1,a)\subseteq A,
\qquad f\longmapsto f(a),
\]
将坐标函数 \(z\mapsto z\) 送到 \(a\)。这里 \(C^*(1,a)\) 是由 \(1,a,a^*\) 生成的范数闭子代数。
\symidx{\(C^*(1,a)\)：由单位元和元素生成的范数闭对合子代数}
\end{theorem}

\begin{proof}
令 \(B=C^*(1,a)\)。正规性使 \(a,a^*\) 交换，故 \(B\) 交换。先核对 \(B\) 内的谱与 \(A\) 内相同。一般地，若自伴元 \(h\in B\) 在 \(A\) 中可逆，取 \(M>\|h\|\)。自伴谱为实数，且不含零，因此正规元的范数公式和多项式谱映射给出
\[
\|1-h^2/M^2\|<1,\qquad
h^{-1}=\frac h{M^2}\sum_{j\ge0}(1-h^2/M^2)^j\in B.
\]
上述级数在 \(B\) 的同一范数下收敛。若一般 \(b\in B\) 在 \(A\) 中可逆，\(b^*b\) 和 \(bb^*\) 都自伴且可逆，故其逆属于 \(B\)。于是 \((b^*b)^{-1}b^*\) 和 \(b^*(bb^*)^{-1}\) 分别为 \(B\) 中的左逆、右逆，二者相同，说明 \(b^{-1}\in B\)。应用于 \(b=a-\lambda1\)，即得谱不变性。

由\cref{b2:thm:gelfand-duality-external}，\(B\cong C(\Delta(B))\)。映射 \(\chi\mapsto\chi(a)\) 的像恰为 \(\sigma_B(a)=\sigma_A(a)\)。它也单射：特征标在 \(a\) 上的值决定在 \(a^*\)、多项式以及其范数极限上的值，所以决定整个 \(B\) 上的特征标。它因而是紧 Hausdorff 空间之间的同胚。把这个同胚与 Gelfand 同构复合，再取逆，就得到所述函数演算，并把坐标函数送到 \(a\)。多项式 \(p(z,\bar z)\) 在谱上稠密，保幺 \(*\)-同态在这些多项式上的值已被 \(a\) 决定，连续性因此给出唯一性。
\end{proof}

这里的\term[lianxu-hanshu-yansuan]{连续函数演算}[continuous functional calculus]允许平方根等连续函数。在上述 \(h(t)=t\) 的例子中，\(h\) 的谱为 \([0,1]\)：若 \(\lambda\) 在区间外，\(1/(t-\lambda)\) 是连续逆函数；若在区间内，\(h-\lambda1\) 在 \(t=\lambda\) 处取零，故不可逆。于是取 \(f(t)=\sqrt t\) 正好得到 \(f(h)=r\)。定理要求元素正规；证明中的谱不变性与函数演算亦见\cite[\nopp II.1.6.7; Corollary II.2.3.1]{Blackadar2006OperatorAlgebras}。

\ProseParagraphBreak
例如 \(a=\operatorname{diag}(1,4)\) 时，谱为 \(\{1,4\}\)，\(f(a)=\operatorname{diag}\bigl(f(1),f(4)\bigr)\)，正平方根为 \(\operatorname{diag}(1,2)\)。相反，\(N=\begin{psmallmatrix}0&1\\0&0\end{psmallmatrix}\) 的谱只有零，而 \(N\ne0\)；若照搬“谱上坐标函数变成 \(N\)”的等距同构，零函数便会变成非零矩阵，矛盾。

\subsection{Gelfand 谱与素谱}
\label{b2:subsec:D02:03}

交换环的素谱由全部素理想组成，Zariski 闭集由理想的共同包含关系给出。Gelfand 谱则由取值同态组成，讨论的是紧或局部紧 Hausdorff 拓扑。特征标的核是闭极大理想，但不能因此把“全部素理想”从定义中删成“特征标的核”，也不能遗漏原来的范数与对合。

\ProseParagraphBreak
有限集合情形两种图景恰好一致：\(\mathbb C^r\) 的素理想都是删去一个坐标的核，得到 \(r\) 个离散点。然而在多项式环 \(\mathbb C[t]\) 中，零理想也是素理想，却不是到 \(\mathbb C\) 的非零取值同态的核。若把多项式限制在 \([0,1]\) 并取一致范数闭包，由\term[Stone-Weierstrass-dingli]{Stone--Weierstrass 定理}[Stone--Weierstrass theorem]得到 \(C\bigl([0,1]\bigr)\)，其 Gelfand 谱是紧区间。这里使用的连续函数逼近定理属于分析学预备知识。这个例子区分了代数素理想的问题与选定范数后完备化的问题；一般素谱留在第三卷第1章，不能用本节的对偶取代。

% !TeX root = ../../Book2.tex
% 纲要标题：### D.3 正性、GNS 与 Gelfand--Naimark 表示
% 纲要位置：工作记录/计划纲要.md，第二卷附录 D。

% 纲要要点（供目录核对）：
% * 正元素、正线性泛函和态（state）；归一化矩阵迹态的具体计算
% * 从态构造内积商空间并完备化，得到 Hilbert 空间上的循环 *-表示
% * 分别证明正锥、平方正性和次序性质，构造达到范数的态；证明 GNS、含幺和非含幺的 Gelfand--Naimark 表示定理及一般 \(*\)-同态的收缩性与单射等距性

\section{正性、GNS 与 Gelfand--Naimark 表示}
\label{b2:sec:D03}

线性泛函可以给代数元素赋予数值，却不必尊重乘法。要求它对平方取非负值，就能够把这些数值组织成内积，再从左乘构造表示。这是本节的主要构造。为让循环向量直接由单位元给出，本节先固定非零含幺 \(C^*\)-代数 \(A\)。

\subsection{正元素、正线性泛函与态}
\label{b2:subsec:D03:01}

\begin{definition}\label{b2:def:positive-element-state}
设 \(A\ne0\) 为含幺复 \(C^*\)-代数，\(a\in A\)。若 \(a=a^*\) 且其谱 \(\sigma_A(a)\subseteq[0,\infty)\)，则称 \(a\) 为\term[zheng-yuansu]{正元素}[positive element]，记 \(a\geq0\)，全体正元素组成 \(A_+\)。若复线性泛函 \(\varphi:A\to\mathbb C\) 对每个 \(a\in A_+\) 都满足 \(\varphi(a)\geq0\)，则称它为\term[zheng-xianxing-fanhan]{正线性泛函}[positive linear functional]。若正线性泛函 \(\varphi\) 还满足 \(\varphi(1_A)=1\)，则称它为 \(A\) 上的\term[tai]{态}[state]。
\symidx{\(A_+\)：\(C^*\)-代数的正元素锥}
\end{definition}

\begin{proposition}[正锥与连续函数演算]\label{b2:prop:cstar-positive-cone}
设 \(A\ne0\) 为含幺复 \(C^*\)-代数，\(A_+\) 为其正元素集合。则 \(A_+\) 是闭凸锥，且 \(A_+\cap(-A_+)=\{0\}\)。每个 \(p\in A_+\) 都有唯一正平方根 \(p^{1/2}\)。每个自伴元素 \(c\in A\) 都可写成
\[
c=c_+-c_-,\qquad c_\pm\in A_+,\quad c_+c_-=0,
\]
其中 \(c_\pm\) 由连续函数 \(t\mapsto\max(\pm t,0)\) 对 \(c\) 作函数演算得到。
\end{proposition}
\begin{proof}
对自伴元素 \(h\)，\cref{b2:thm:continuous-functional-calculus}把 \(C^*(1,h)\) 识别为实谱上的连续函数代数。因此 \(h\ge0\) 等价于谱上坐标函数非负，且
\[
h\ge0,\ \|h\|\le1
\quad\Longleftrightarrow\quad
h=h^*,\ \|h\|\le1,\ \|1-h\|\le1.
\]
右侧是三个闭凸集的交，故 \(A_+\) 与单位球的交闭且凸。将有界收敛序列或有限多个正元素同时缩放到单位球中，就得到正锥的闭性与凸性。非负实数倍仍为正，故它是闭凸锥。若 \(h\) 与 \(-h\) 都正，其谱只能是 \(\{0\}\)，自伴元素的范数等于谱半径，故 \(h=0\)。

对正元素的谱使用 \(t\mapsto\sqrt t\)，同一函数演算给出 \(p^{1/2}\)。对任意自伴元素 \(c\)，使用 \(t\mapsto\max(t,0)\) 和 \(t\mapsto\max(-t,0)\)，便得到所述分解及 \(c_+c_-=0\)。若正元素 \(q\) 满足 \(q^2=p\)，则 \(q\) 与 \(p\) 交换，也与 \(p^{1/2}\) 交换。在这两个元素生成的交换 \(C^*\)-代数中，非负函数的平方根逐点唯一，故 \(q=p^{1/2}\)。
\end{proof}

正平方根把谱上的非负性与代数中的平方联系起来。还须证明：即使 \(x\) 不是正规元素，\(x^*x\) 仍落在正锥中。

\begin{lemma}[平方元素的正性]\label{b2:lem:cstar-square-positive}
设 \(A\ne0\) 为含幺复 \(C^*\)-代数，则对每个 \(x\in A\)，都有 \(x^*x\in A_+\)。
\end{lemma}
\begin{proof}
对任意 \(u,v\in A\)，\(1-uv\) 可逆当且仅当 \(1-vu\) 可逆，因为一个方向的逆为
\[
(1-vu)^{-1}=1+v(1-uv)^{-1}u,
\]
另一方向交换 \(u,v\) 即得。对 \(\lambda\ne0\)，取 \(u=\lambda^{-1}y^*\)、\(v=y\)，便得
\[
\sigma_A(y^*y)\setminus\{0\}
=\sigma_A(yy^*)\setminus\{0\}.
\]
两者都是自伴元素，而可能多出的谱点 \(0\) 也属于 \([0,\infty)\)，所以其中一个为正当且仅当另一个为正。

写 \(x=a+ib\)，其中 \(a,b\) 自伴。由函数演算，\(a^2,b^2\ge0\)，再由正锥的加法封闭性得到 \(x^*x+xx^*=2a^2+2b^2\ge0\)。令 \(c=x^*x=c_+-c_-\)，取 \(y=xc_-\)。于是
\[
y^*y=c_-cc_-=-c_-^3,\qquad
yy^*=(y^*y+yy^*)+c_-^3\ge0.
\]
第二式对 \(y\) 使用同样的平方和等式，并使用 \(c_-^3\ge0\)。非零谱相同又使 \(y^*y\ge0\)；它的相反数也正，故 \(c_-^3=0\)。正元素 \(c_-\) 的谱因而只能是 \(\{0\}\)，于是 \(c_-=0\)。因此 \(x^*x\ge0\)。
\end{proof}

\begin{proposition}[正元素的次序性质]\label{b2:prop:cstar-positive-input}
设 \(A\ne0\) 为含幺复 \(C^*\)-代数。对自伴元素 \(s,t\in A\)，以 \(t-s\ge0\) 定义 \(s\le t\)。这给出自伴部分上的偏序；对任意 \(a,b\in A\)、正元素 \(p\in A_+\) 及自伴元素 \(h\in A\)，有
\[
0\le a^*a\le\|a\|^2 1,\qquad
b^*pb\ge0,\qquad
-\|h\|1\le h\le\|h\|1.
\]
\end{proposition}
\begin{proof}
正锥的加法封闭性和 \(A_+\cap(-A_+)=\{0\}\) 分别给出传递性和反对称性，\(0\in A_+\) 给出自反性。由上一个引理及 \(C^*\) 恒等式，\(a^*a\) 的谱包含于 \([0,\|a^*a\|]=[0,\|a\|^2]\)，故函数演算给出第一组界。又有
\[
b^*pb=(p^{1/2}b)^*(p^{1/2}b)\ge0.
\]
自伴元素 \(h\) 的谱位于 \([-\|h\|,\|h\|]\)，对 \(t\mapsto\|h\|\pm t\) 作函数演算便得到最后一组界。
\end{proof}

这些谱论与次序性质也见 Blackadar\cite[Propositions II.3.1.2--II.3.1.3; Corollary II.3.1.5; II.3.1.8]{Blackadar2017OperatorAlgebrasCorrected}。矩阵情形的正元素恰为半正定 Hermitian 矩阵，函数情形则是逐点非负的连续函数。

\begin{lemma}\label{b2:lem:positive-functional-cs}
设 \(A\ne0\) 为含幺复 \(C^*\)-代数，\(\varphi:A\rightarrow\mathbb C\) 为正线性泛函，则对任意 \(a,b\in A\)，有
\[
\varphi(a^*)=\overline{\varphi(a)},\qquad
\bigl|\varphi(a^*b)\bigr|^2\le\varphi(a^*a)\varphi(b^*b).
\]
它自动有界，且 \(\|\varphi\|=\varphi(1)\)。
\end{lemma}
\begin{proof}
对自伴 \(h\)，\cref{b2:prop:cstar-positive-input}使 \(\|h\|1\pm h\) 正，故 \(\varphi(h)\) 为实数。将任意 \(a\) 写成 \(h+ik\)、\(h,k\) 自伴，即得第一式。令 \([a,b]=\varphi(a^*b)\)，这是对第一变量共轭线性、第二变量线性的半正定 Hermitian 形式。

若 \([b,b]>0\)，代入 \(a-\lambda b\)、\(\lambda=[b,a]/[b,b]\)，由其平方值非负得到 \(\bigl|[a,b]\bigr|^2\le[a,a][b,b]\)。若 \([b,b]=0\)，非负式 \([a+zb,a+zb]\ge0\) 对任意复数 \(z\) 成立，其关于 \(z\) 的线性项只能为零，故 \([a,b]=0\)。这证明两种情形。

再把不等式用于 \(a,1\)，由 \(\varphi(a^*)=\overline{\varphi(a)}\) 及 \(0\le a^*a\le\|a\|^2 1\) 得
\[
\bigl|\varphi(a)\bigr|^2\le\varphi(a^*a)\varphi(1)
\le\|a\|^2\cdot\varphi(1)^2.
\]
所以 \(\|\varphi\|\le\varphi(1)\)，而在单位元处取值给出反向不等式。
\end{proof}

例如 \(M_n(\mathbb C)\) 上的归一化迹 \(\tau(a)=\operatorname{tr}(a)/n\) 是态，因为 \(\operatorname{tr}(a^*a)=\sum_{i,j}|a_{ij}|^2\)。对单位向量 \(v\)，\(\omega_v(a)=\langle v,av\rangle\) 也是态；\(C(X)\) 中的点取值则既是态又保持乘法。一般态不保持乘法，如 \(M_2(\mathbb C)\) 中 \(\tau(E_{11})^2=1/4\ne1/2=\tau(E_{11}^2)\)。

\ProseParagraphBreak
态却能提供长度和内积的候选：\(\varphi(a^*b)\) 对第一变量共轭线性、对第二变量线性，且平方值非负。这个型可能退化。例如取 \(v=e_1\)，则 \(\omega_v(a^*a)=\|ae_1\|^2\)；非零矩阵 \(E_{22}\) 的值为零，因为 \(E_{22}e_1=0\)。所以要从态得到真正的内积空间，必须先把所有零长度元素模去，再作完备化。下面的 GNS 构造正是将这两步与代数的左乘作用结合起来。

\subsection{GNS 构造与循环表示}
\label{b2:subsec:D03:02}

\begin{theorem}[Gelfand--Naimark--Segal 构造]\label{b2:thm:gns-unital}
设 \(A\ne0\) 为含幺复 \(C^*\)-代数，\(\varphi\) 为 \(A\) 上的态，则存在复 Hilbert 空间 \(H_\varphi\)、保幺 \(*\)-表示 \(\pi_\varphi:A\to B(H_\varphi)\) 和单位向量 \(\xi_\varphi\)，使对每个 \(a\in A\)，
\[
\varphi(a)=\bigl\langle\xi_\varphi,\pi_\varphi(a)\xi_\varphi\bigr\rangle,
\qquad\overline{\pi_\varphi(A)\xi_\varphi}=H_\varphi.
\]
这个三元组在保持指定向量并交织表示的酉同构下唯一。
\symidx{\(H_\varphi\)：态的 GNS Hilbert 空间}
\symidx{\(\pi_\varphi\)：态的 GNS 表示}
\symidx{\(\xi_\varphi\)：GNS 表示的指定循环向量}
\end{theorem}
\begin{proof}
令
\[
N_\varphi=\bigl\{\,a\in A:\varphi(a^*a)=0\,\bigr\}.
\]
\symidx{\(N_\varphi\)：GNS 构造中半内积的零空间}
由\cref{b2:lem:positive-functional-cs}，其中每个元素与任意元素的半内积都为零，故 \(N_\varphi\) 为线性子空间，且
\[
\bigl\langle[a],[b]\bigr\rangle=\varphi(a^*b)
\]
在商 \(A/N_\varphi\) 上良定义并正定。\cref{b2:prop:cstar-positive-input}还给出
\begin{equation}\label{b2:eq:gns-bounded-left-action}
\varphi\bigl((ca)^*(ca)\bigr)
=\varphi(a^*c^*ca)\le\|c\|^2\cdot\varphi(a^*a).
\end{equation}
确实，对正元素 \(\|c\|^2 1-c^*c\) 作 \(a^*(\cdot)a\) 并应用 \(\varphi\)，就得到此不等式。它同时证明 \(N_\varphi\) 为左理想，以及左乘 \([a]\mapsto[ca]\) 为范数至多 \(\|c\|\) 的有界算子。

把内积空间完备化为 \(H_\varphi\)，左乘唯一连续延拓为 \(\pi_\varphi(c)\)。乘法及单位元关系先在稠密商空间上成立，连续性使其在完备化后仍成立。又有
\[
\bigl\langle[ca],[b]\bigr\rangle=\varphi(a^*c^*b)
=\bigl\langle[a],[c^*b]\bigr\rangle,
\]
所以 \(\pi_\varphi(c)^*=\pi_\varphi(c^*)\)。取 \(\xi_\varphi=[1]\)，它的范数为一，\(\pi_\varphi(a)\xi_\varphi=[a]\)，因而满足循环性及泛函公式。

若 \((\rho,H,\eta)\) 是另一三元组，在稠密子空间上定义 \([a]\mapsto\rho(a)\eta\)。其内积为
\[
\bigl\langle\rho(a)\eta,\rho(b)\eta\bigr\rangle=\varphi(a^*b),
\]
故良定义且等距。循环性使其像稠密，所以延拓为满酉算子；它交织左作用并送指定向量到 \(\eta\)。这两个要求又在稠密向量 \([a]\) 上决定它，故唯一。
\end{proof}

这称为\term[GNS-gouzao]{GNS 构造}[Gelfand--Naimark--Segal construction]。\footnote{奈马克（\OriginalName{Марк Аронович Наймарк}，1909-1978），苏联数学家，研究泛函分析与算子表示。西格尔（Irving Ezra Segal，1918-1998），美国数学家，研究泛函分析与量子场论。}上面的完备化是必需的一步，不能把一般的代数商空间直接当作 Hilbert 空间。非含幺版本及其循环向量的构造可参阅\cite[\nopp II.6.4.1--II.6.4.3]{Blackadar2017OperatorAlgebrasCorrected}。

\ProseParagraphBreak
对 \(M_n(\mathbb C)\) 的迹态，\(N_\tau=0\)，所以 \(H_\tau=M_n(\mathbb C)\)，内积为 \(\operatorname{tr}(a^*b)/n\)，表示为左乘，循环向量为 \(I_n\)。维数为 \(n^2\)，并非自然列模的 \(n\)。对向量态 \(\omega_v\)，则 \(N_{\omega_v}=\{\,a:av=0\,\}\)，映射 \([a]\mapsto av\) 等距满射到 \(\mathbb C^n\)，恢复自然表示。不同态因此可以给出不同的循环模型。

\subsection{Gelfand--Naimark 表示定理}
\label{b2:subsec:D03:03}

一次 GNS 构造未必忠实。例如 \(C(X)\) 的点取值只看到一个点，所有在该点为零的函数都被零化。为了看见整个代数，需要足够多的态。

\begin{proposition}[达到范数的态]\label{b2:prop:norming-state-external}
设 \(A\ne0\) 为含幺复 \(C^*\)-代数。对每个正元素 \(p\in A_+\)，存在 \(A\) 上的态 \(\varphi\) 使 \(\varphi(p)=\|p\|\)。
\end{proposition}

\begin{proof}
由函数演算，\(\|p\|\) 属于 \(p\) 的谱；在 \(C^*(1,p)\cong C(\sigma_A(p))\) 上取该点的取值泛函 \(\psi\)。它满足 \(\|\psi\|=\psi(1)=1\)、\(\psi(p)=\|p\|\)。\term[Hahn-Banach-dingli]{Hahn--Banach 定理}[Hahn--Banach theorem]将它保范延拓为 \(A\) 上的复线性泛函 \(\varphi\)，仍有 \(\|\varphi\|=\varphi(1)=1\)。

核对延拓的正性。对自伴元 \(h\) 及实数 \(t\)，\(C^*\) 恒等式给出
\[
|1+it\cdot\varphi(h)|^2\le\|1+ith\|^2
=\|1+t^2h^2\|\le1+t^2\|h\|^2.
\]
左端关于 \(t\) 的一次项为 \(-2t\operatorname{Im}\varphi(h)\)。让 \(t\) 从正、负两侧趋于零，必得 \(\operatorname{Im}\varphi(h)=0\)。若 \(q\ge0\)、\(\|q\|\le1\)，正性的谱论基础使 \(\|1-q\|\le1\)，故 \(|1-\varphi(q)|\le1\)。因为 \(\varphi(q)\) 为实数，必有 \(\varphi(q)\ge0\)。任意正元素可缩放到这种情形，所以 \(\varphi\) 是态，并保留所需取值。
\end{proof}

这里仅用 Hahn--Banach 的保范延拓作为分析学基础；延拓何以仍为正已在证明中核对。这个存在性论证亦见\cite[\nopp II.6.3.1; Proposition II.6.3.3]{Blackadar2017OperatorAlgebrasCorrected}。

\begin{theorem}[Gelfand--Naimark，含幺情形]\label{b2:thm:gelfand-naimark-unital}
设 \(A\ne0\) 为含幺复 \(C^*\)-代数，则存在复 Hilbert 空间 \(H\) 及保幺等距 \(*\)-同态 \(A\rightarrow B(H)\)，其像为 \(B(H)\) 的含幺范数闭子代数。
\end{theorem}
\begin{proof}
记全部态组成的集合为 \(\mathcal S(A)\)。先取有限支集的代数直和 \(\bigoplus_{\varphi\in\mathcal S(A)}H_\varphi\)，以各分量内积之和为内积，再作完备化，得到 Hilbert 直和 \(H\)。因为各 \(H_\varphi\) 已经完备，\(H\) 可识别为
\[
\begin{gathered}
H=\left\{\,(\xi_\varphi)_{\varphi\in\mathcal S(A)}
\,\Bigm\vert\,
\sum_{\varphi\in\mathcal S(A)}\|\xi_\varphi\|^2<\infty\,\right\},
\\
\|\xi\|^2=\sum_\varphi\|\xi_\varphi\|^2.
\end{gathered}
\]
任意指标集上的平方和指有限子集部分和的上确界。因此这里允许平方可和的分量族，而不仅是有限支集向量。

令 \(\pi(a)(\xi_\varphi)_\varphi=(\pi_\varphi(a)\xi_\varphi)_\varphi\)。式\eqref{b2:eq:gns-bounded-left-action}给出与态无关的统一界 \(\|\pi_\varphi(a)\|\le\|a\|\)，于是
\[
\sum_\varphi\|\pi_\varphi(a)\xi_\varphi\|^2
\le\|a\|^2\sum_\varphi\|\xi_\varphi\|^2.
\]
故 \(\pi(a)\) 确实把 \(H\) 送入自身，且 \(\|\pi(a)\|\le\|a\|\)。乘法、单位与对合关系先在有限支集的稠密子空间上成立，再由连续性延伸到 \(H\)，所以 \(\pi\) 是含幺 \(*\)-同态。对给定 \(a\)，应用\cref{b2:prop:norming-state-external}于 \(p=a^*a\)，选择相应态，把其单位循环向量 \(\xi_\varphi\) 放在 \(H\) 的该分量、其他分量置零，则
\[
\bigl\lVert\pi(a)\bigr\rVert^2\ge\bigl\lVert\pi_\varphi(a)\xi_\varphi\bigr\rVert^2
=\varphi(a^*a)=\|a^*a\|=\|a\|^2.
\]
结合相反方向的统一界得到等距性，特别是单射。等距像因 \(A\) 完备而闭，结论成立。
\end{proof}

这就是含幺情形的\term[Gelfand-Naimark-biaoshi-dingli]{Gelfand--Naimark 表示定理}[Gelfand--Naimark representation theorem]。本证明先构造态，再以它们的 GNS 表示组成忠实表示。

\begin{proposition}[添幺与非含幺表示]\label{b2:prop:cstar-unitization-representation}
设 \(A\ne0\) 为不含幺的复 \(C^*\)-代数。它可等距嵌入一个含幺 \(C^*\)-代数 \(\widetilde A\)，作为闭双边 \(*\)-理想，且 \(\widetilde A/A\cong\mathbb C\)。特别地，\(A\) 有忠实等距 Hilbert 空间表示。
\end{proposition}
\begin{proof}
取向量空间 \(\widetilde A=A\oplus\mathbb C1\)，规定
\[
(a+\lambda1)(b+\mu1)=ab+\lambda b+\mu a+\lambda\mu1,
\qquad (a+\lambda1)^*=a^*+\bar\lambda1.
\]
把它作用于 Banach 空间 \(A\)，令 \(L_{a+\lambda1}x=ax+\lambda x\)，定义 \(N(c)=\|L_c\|\)。若 \(N(a+\lambda1)=0\) 且 \(\lambda\ne0\)，\(-a/\lambda\) 就是 \(A\) 的左单位；取对合得到右单位，两者相同，违反不含幺假设。若 \(\lambda=0\)，则 \(aa^*=0\)，故 \(a=0\)。因此 \(N\) 是范数，且由 \(x=a^*/\|a\|\) 得 \(N(a)=\|a\|\)。

对 \(c\in\widetilde A\)、\(x\in A\)，有
\[
\|c^*x\|^2=\|x^*cc^*x\|
\le\|x\|\cdot N(c)\cdot\|c^*x\|.
\]
消去非零因子并取上确界，得 \(N(c^*)\le N(c)\)；取对合又得相等。再由
\[
\|cx\|^2=\|x^*c^*cx\|\le\|x\|^2\cdot N(c^*c)
\]
和次乘性，得到 \(N(c)^2=N(c^*c)\)。

还须检查完备性。必有 \(|\lambda|\le N(a+\lambda1)\)：否则在 \(B(A)\) 中有 \(\|I+L_{a/\lambda}\|<1\)。写成
\[
-L_{a/\lambda}=I-\bigl(I+L_{a/\lambda}\bigr),
\]
右侧的括号内算子范数小于一，故 Neumann 级数给出 \(-L_{a/\lambda}\) 的逆，从而 \(L_a\) 可逆。取 \(e\in A\) 满足 \(ae=a\)，则 \(a(ex-x)=0\) 及 \(L_a\) 单射给出 \(ex=x\)，又会产生单位元。因此 \(\|a\|\le2N(a+\lambda1)\)。一个 \(N\)-柯西列的 \(A\) 坐标和复数坐标分别柯西，取两坐标的极限便得它的 \(N\)-极限。于是 \(\widetilde A\) 确为含幺 \(C^*\)-代数，\(A\) 为其中闭理想，取复数坐标给出商同构。将上一定理的等距表示限制到 \(A\)，即得非含幺结论。
\end{proof}

这个表示的原代数像未必含有 Hilbert 空间的单位算子；添幺和一般表示的比较见\cite[\nopp II.1.2; Corollary II.6.4.10]{Blackadar2017OperatorAlgebrasCorrected}。

\begin{proposition}[一般 \(C^*\)-代数的 \(*\)-同态]\label{b2:prop:cstar-star-hom-general}
设 \(A,B\) 为复 \(C^*\)-代数，\(\pi:A\to B\) 为 \(*\)-同态，不要求代数含幺或映射保幺。则对每个 \(a\in A\)，有
\[
\|\pi(a)\|\le\|a\|.
\]
因此 \(\pi\) 自动连续。若 \(\pi\) 单射，则它等距，即 \(\|\pi(a)\|=\|a\|\)。
\end{proposition}
\begin{proof}
零代数的情形直接成立。若 \(A\) 含幺，结论由\cref{b2:prop:cstar-star-hom-contractivity-unital,b2:prop:cstar-star-hom-isometry-unital}给出。

设 \(A\) 不含幺，取它的等距添幺 \(\widetilde A\)。若 \(B\) 含幺，令 \(B^\sharp=B\)；否则取 \(B^\sharp=\widetilde B\)。公式
\[
\widetilde\pi(a+\lambda1)=\pi(a)+\lambda1_{B^\sharp}
\]
定义保幺 \(*\)-同态 \(\widetilde\pi:\widetilde A\to B^\sharp\)。含幺情形的收缩性及添幺的等距嵌入给出
\(\|\pi(a)\|\le\|a\|\)。

再设 \(\pi\) 单射。若 \(\widetilde\pi(a+\lambda1)=0\)，当 \(\lambda=0\) 时已有 \(a=0\)；若 \(\lambda\ne0\)，则 \(e=-a/\lambda\in A\) 满足 \(\pi(e)=1_{B^\sharp}\)。对每个 \(x\in A\)，于是
\[
\pi(ex)=\pi(x)=\pi(xe).
\]
由单射性得 \(ex=xe=x\)，这与 \(A\) 不含幺矛盾。因此 \(\widetilde\pi\) 也单射，含幺情形的等距性限制到 \(A\) 即得所需结论。
\end{proof}

% !TeX root = ../../Book2.tex
% 纲要标题：### D.4 von Neumann 代数与双交换子

% 纲要要点（供目录核对）：
% * 弱算子拓扑、强算子拓扑与范数拓扑；含恒等算子的 *-子代数的闭包
% * 双交换子定理的含幺自伴版本：\(A''=\overline A^{\mathrm{SOT}}=\overline A^{\mathrm{WOT}}\)
% * 中心、因子、投影与迹；以投影等价、有限性及极小性介绍 I、II、III 型因子及 II_1、II_∞ 的区别，计算 \(B(H)\) 的投影模型；不展开完整分类与 Tomita–Takesaki 模理论
% #### D.4.1 弱算子、强算子与范数拓扑
% #### D.4.2 双交换子定理
% #### D.4.3 中心、因子、投影与迹
% #### D.4.4 投影等价与因子的类型

\section{von Neumann 代数与双交换子}
\label{b2:sec:D04}

当 Hilbert 空间 \(H\) 有限维时，\(B(H)\) 上的范数、强算子与弱算子拓扑一致；无限维时则必须区分怎样取极限。双交换子定理的内容，正是把一种算子拓扑的闭包与代数交换关系联系起来。进入弱闭代数以后，还可以通过其中投影的比较，区分有限维矩阵没有表现出的结构。本节始终令 \(H\ne0\) 为复 Hilbert 空间，所有子代数均作为 \(B(H)\) 的具体子集来讨论。

\subsection{弱算子、强算子与范数拓扑}
\label{b2:subsec:D04:01}

\begin{definition}\label{b2:def:operator-topologies}
设 \(H\) 为复 Hilbert 空间，\((T_\alpha)\) 是 \(B(H)\) 中以有向集为指标的网，\(T\in B(H)\)。若对每个 \(\xi\in H\)，都有 \(\bigl\lVert(T_\alpha-T)\xi\bigr\rVert\to0\)，则称 \((T_\alpha)\) 在\term[qiang-suanzi-tuopu]{强算子拓扑}[strong operator topology]中趋于 \(T\)；若对每对 \(\xi,\eta\in H\)，都有 \(\bigl\langle(T_\alpha-T)\xi,\eta\bigr\rangle\to0\)，则称它在\term[ruo-suanzi-tuopu]{弱算子拓扑}[weak operator topology]中趋于 \(T\)。两者简称 SOT、WOT；范数收敛则要求 \(\|T_\alpha-T\|\to0\)。
\end{definition}
\symidx{SOT：强算子拓扑}
\symidx{WOT：弱算子拓扑}

这里的网是以有向集为指标的族，极限要求从某个指标起均满足给定邻域条件；不能一般只用序列描述闭包。SOT 的基本邻域同时限制有限多个 \(\bigl\lVert(S-T)\xi_j\bigr\rVert\)，WOT 则同时限制有限多个矩阵系数。直接由定义和 Cauchy--Schwarz 不等式可得
\[
\text{范数收敛}\ \Longrightarrow\ \text{SOT 收敛}\ \Longrightarrow\ \text{WOT 收敛}.
\]

在 \(\ell^2(\mathbb N)\) 中，令 \(P_n\) 投影到前 \(n\) 个标准基向量。对 \(\xi=(\xi_j)\)，尾和 \(\bigl\lVert(I-P_n)\xi\bigr\rVert^2=\sum_{j>n}|\xi_j|^2\to0\)，故 \(P_n\to I\) 强收敛；但 \(\|I-P_n\|=1\)，并不范数收敛。再令 \(S e_j=e_{j+1}\)。对有限支集向量 \(\xi,\eta\)，\(\langle S^n\xi,\eta\rangle\) 最终为零；利用有限支集向量的稠密性和 \(\|S^n\|=1\)，可推广到任意 \(\xi,\eta\)。因此 \(S^n\to0\) 弱收敛，但 \(\|S^n e_1\|=1\)，不强收敛。

\subsection{双交换子定理}
\label{b2:subsec:D04:02}

设 \(H\) 为复 Hilbert 空间，\(\mathcal S\subseteq B(H)\)。满足与 \(\mathcal S\) 中每个算子交换的全体有界算子组成\term[jiaohuanzi-daishu]{交换子代数}[commutant]
\[
\mathcal S'=\bigl\{\,T\in B(H):TS=ST\text{ 对所有 }S\in\mathcal S\,\bigr\},
\qquad \mathcal S''=(\mathcal S')'.
\]
\symidx{\(\mathcal S'\)：有界算子集合的交换子代数}
\symidx{\(\mathcal S''\)：有界算子集合的双交换子}
这与第十一章的中心化子是同类代数条件，但环境现在是全部有界算子，闭包也有指定的拓扑。

\begin{theorem}[von Neumann 双交换子]\label{b2:thm:operator-bicommutant}
设 \(H\) 为复 Hilbert 空间，\(\mathcal A\subseteq B(H)\) 是含恒等算子 \(I_H\) 的 \(*\)-子代数。记 \(\mathcal A'=\{T\in B(H):Ta=aT\text{ 对所有 }a\in\mathcal A\}\)，\(\mathcal A''=(\mathcal A')'\)。下式中的上标 SOT 和 WOT 分别表示强算子拓扑与弱算子拓扑下的闭包。此时
\[
\mathcal A''=\overline{\mathcal A}^{\mathrm{SOT}}
=\overline{\mathcal A}^{\mathrm{WOT}}.
\]
\end{theorem}
\begin{proof}
固定 \(S\)，等式 \(TS=ST\) 在 WOT 下是闭条件，因为其每个矩阵系数是
\(\langle TS\xi,\eta\rangle-\langle T\xi,S^*\eta\rangle\)，对变量 \(T\) 连续。所以交换子代数为 WOT 闭集。又 \(\mathcal A\subseteq\mathcal A''\)，故
\[
\overline{\mathcal A}^{\mathrm{SOT}}\subseteq
\overline{\mathcal A}^{\mathrm{WOT}}\subseteq\mathcal A''.
\]
只需证明 \(T\in\mathcal A''\) 可以在任意有限组向量上同时由 \(\mathcal A\) 逼近。

给定 \(\xi_1,\ldots,\xi_m\)，在 \(H^m\) 中令
\[
\mathcal M=\overline{\bigl\{\,(a\xi_1,\ldots,a\xi_m):a\in\mathcal A\,\bigr\}},
\]
闭包使用 Hilbert 范数。这是闭线性子空间，对每个对角作用 \(a^{(m)}\) 不变；由于 \(a^*\in\mathcal A\)，它的正交补也不变。故正交投影 \(P\) 与每个 \(a^{(m)}\) 交换。把 \(P\) 按 \(H^m\) 的分量写成算子矩阵 \((P_{ij})\)，交换关系逐条目给出 \(P_{ij}a=aP_{ij}\)，即 \(P_{ij}\in\mathcal A'\)。

因为 \(T\in\mathcal A''\)，它与每个 \(P_{ij}\) 交换，从而 \(T^{(m)}P=PT^{(m)}\)。单位元使 \(\xi=(\xi_1,\ldots,\xi_m)\in\mathcal M\)，所以
\(T^{(m)}\xi=T^{(m)}P\xi=PT^{(m)}\xi\in\mathcal M\)。依照 \(\mathcal M\) 的定义，对任意 \(\varepsilon>0\)，存在同一个 \(a\in\mathcal A\) 使
\[
\sum_{j=1}^m\bigl\lVert(T-a)\xi_j\bigr\rVert^2<\varepsilon^2.
\]
这就使 \(a\) 落在指定 SOT 邻域中，证明 \(T\in\overline{\mathcal A}^{\mathrm{SOT}}\)。
\end{proof}

这称为\term[von-Neumann-shuangjiaohuanzi-dingli]{von Neumann 双交换子定理}[von Neumann bicommutant theorem]。\footnote{冯·诺伊曼（Neumann János Lajos，1903-1957），匈牙利裔美国数学家，奠定算子代数理论，并研究量子力学与计算机结构。}论证中，对伴随封闭用于让正交补不变，单位元用于确保 \(\xi\) 落在循环子空间，有限次放大则用于把“一次逼近一个向量”加强到“同一个算子同时逼近有限多个向量”。可对照 Blackadar\cite[Theorem I.9.1.1; Corollary I.9.1.2]{Blackadar2017OperatorAlgebrasCorrected}，其中还给出非退化作用的版本；此处只使用已经证明的含幺版本。

\begin{definition}\label{b2:def:von-neumann-algebra}
设 \(H\) 为复 Hilbert 空间。若 \(\mathcal M\subseteq B(H)\) 是含恒等算子 \(I_H\) 且在弱算子拓扑下闭的 \(*\)-子代数，则称 \(\mathcal M\) 为\term[von-Neumann-daishu]{von Neumann 代数}[von Neumann algebra]。这里 \(\mathcal M'\) 是 \(\mathcal M\) 在 \(B(H)\) 中的交换子代数，\(\mathcal M''=(\mathcal M')'\)。由\cref{b2:thm:operator-bicommutant}，这等价于 \(\mathcal M=\mathcal M''\)，也等价于含幺 \(*\)-子代数在强算子拓扑下闭。
\end{definition}

\subsection{中心、因子、投影与迹}
\label{b2:subsec:D04:03}

设 \(H\) 为复 Hilbert 空间，\(\mathcal M\subseteq B(H)\) 为 von Neumann 代数。其中心为 \(Z(\mathcal M)=\mathcal M\cap\mathcal M'\)。如果 \(Z(\mathcal M)=\mathbb C I_H\)，则称 \(\mathcal M\) 为\term[yinzi-suanzi]{因子}[factor]。若 \(p\in\mathcal M\) 满足 \(p=p^*=p^2\)，则称 \(p\) 为正交投影；它对应 \(H\) 的一个闭子空间，但并非 \(H\) 的每个闭子空间的投影都属于 \(\mathcal M\)。设 \(\tau\) 是 \(\mathcal M\) 上的态。若对任意 \(a,b\in\mathcal M\) 都有 \(\tau(ab)=\tau(ba)\)，则称 \(\tau\) 为\term[jiantai]{迹态}[tracial state]；若对每个非零 \(a\in\mathcal M\) 还有 \(\tau(a^*a)>0\)，则称它为忠实迹态。

\ProseParagraphBreak
\(B(H)\) 是因子。若 \(T\) 与全部秩一正交投影交换，它保持每条直线，故 \(T\xi=\lambda_\xi\xi\)。对两个线性无关向量再考虑它们的和，得到相同的标量，于是 \(T=\lambda I\)。有限维时 \(M_n(\mathbb C)\) 的归一化迹为忠实迹态；而 \(\bigoplus_{j=1}^rM_{n_j}(\mathbb C)\) 的中心为 \(\mathbb C^r\)，\(r>1\) 时不是因子，其迹态包括
\[
\tau(a_1,\ldots,a_r)=\sum_{j=1}^r t_j\frac{\operatorname{tr}(a_j)}{n_j},
\qquad t_j\ge0,\quad\sum_jt_j=1.
\]
所有 \(t_j>0\) 时该态忠实。无限维时，投影像空间的维数已不能单独刻画它在代数中的大小：比较所用的算子是否属于当前代数，同样是条件的一部分。

\subsection{投影等价与因子的类型}
\label{b2:subsec:D04:04}

一个闭子空间可以与自己的真子空间等距同构，例如单边移位把 \(\ell^2(\mathbb N)\) 等距映入去掉第一个坐标的子空间。投影的有限性要问的是：这种等距映射能否由当前 von Neumann 代数中的算子实现。

\begin{definition}\label{b2:def:projection-equivalence}
设 \(H\) 为复 Hilbert 空间，\(\mathcal M\subseteq B(H)\) 为 von Neumann 代数。若 \(v\in\mathcal M\) 满足 \(vv^*v=v\)，则称 \(v\) 为\term[bufen-dengju]{部分等距算子}[partial isometry]。此时 \(v^*v\) 与 \(vv^*\) 都是投影，分别称为它的初投影与终投影。对投影 \(p,q\in\mathcal M\)，若存在部分等距算子 \(v\in\mathcal M\)，使
\[
v^*v=p,\qquad vv^*=q,
\]
则称 \(p\) 与 \(q\) 在 \(\mathcal M\) 中等价，记为 \(p\sim_{\mathcal M}q\)。这一关系称为\term[Murray-von-Neumann-dengjia]{Murray--von Neumann 投影等价}[Murray--von Neumann equivalence]。\footnote{默里（Francis Joseph Murray，1911-1996），美国数学家，与 von Neumann 合作奠定算子环及因子理论。}
\end{definition}

由 \(vv^*v=v\) 可直接验证两个自伴算子的平方仍为自身。算子 \(v\) 在 \(pH\) 上保持范数，像恰为 \(qH\)，在 \((1-p)H\) 上为零。取 \(v=p\) 给出自反性，取伴随给出对称性；若 \(v\) 把 \(p\) 连接到 \(q\)，而 \(w\) 把 \(q\) 连接到 \(r\)，则 \((wv)^*(wv)=p\)、\((wv)(wv)^*=r\)，得到传递性。下文写 \(p\sim q\) 时，等价都在指定的 \(\mathcal M\) 内。

\begin{definition}\label{b2:def:finite-minimal-projection}
设 \(H\) 为复 Hilbert 空间，\(\mathcal M\subseteq B(H)\) 为 von Neumann 代数，\(p\in\mathcal M\) 为投影。对另一投影 \(q\in\mathcal M\)，规定 \(q\le p\) 表示 \(pq=qp=q\)，即 \(qH\subseteq pH\)。若对每个满足 \(q\le p\) 且 \(q\sim_{\mathcal M}p\) 的投影 \(q\)，都有 \(q=p\)，则称 \(p\) 为\term[youxian-touying]{有限投影}[finite projection]；否则称它为\term[wuxian-touying]{无限投影}[infinite projection]。若 \(p\ne0\)，且它的子投影只有 \(0\) 和 \(p\)，则称 \(p\) 为\term[jixiao-touying]{极小投影}[minimal projection]。
\end{definition}

极小投影必为有限投影，因为与非零投影等价的投影仍非零。“有限”取决于环境代数；例如在无限维 \(H\) 上，\(I_H\) 在 \(\mathbb C I_H\) 内为极小投影，却在 \(B(H)\) 内为无限投影。

\begin{proposition}\label{b2:prop:finite-projection-models}
设 \(H\ne0\) 为复 Hilbert 空间。\(B(H)\) 中的投影有限，当且仅当它的像空间有限维；非零投影极小，当且仅当它的像空间一维。若 von Neumann 代数 \(\mathcal M\subseteq B(H)\) 有忠实迹态，则它的每个投影均有限。
\end{proposition}
\begin{proof}
在 \(B(H)\) 中，\(p\sim q\) 正好表示 \(pH\) 与 \(qH\) 酉同构：一个方向由部分等距算子的限制给出，另一个方向把酉同构在 \((1-p)H\) 上延拓为零。有限维空间不能与自己的真子空间酉同构，所以有限秩投影有限。

若 \(pH\) 无限维，在其中取正交单位向量 \(\xi_1,\xi_2,\ldots\)。令 \(v\xi_j=\xi_{j+1}\)，在这列向量闭线性张成空间于 \(pH\) 内的正交补上令 \(v\) 为恒等，在 \((1-p)H\) 上令它为零。于是 \(v^*v=p\)，而 \(vv^*=p-e\)，其中 \(e\) 为到 \(\mathbb C\xi_1\) 的投影。故 \(p\) 为无限投影。一维空间没有非零真子空间，维数大于一时则有，所以极小性也得到刻画。

最后设 \(\tau\) 为 \(\mathcal M\) 上的忠实迹态，\(q\le p\) 且 \(v^*v=p\)、\(vv^*=q\)。迹性质给出 \(\tau(p)=\tau(q)\)，从而 \(\tau(p-q)=0\)。由于 \(p-q\) 为投影，忠实性迫使 \(p-q=0\)。
\end{proof}

\begin{definition}\label{b2:def:factor-types}
设 \(H\ne0\) 为复 Hilbert 空间，\(\mathcal M\subseteq B(H)\) 为 von Neumann 代数，并且 \(Z(\mathcal M)=\mathbb C I_H\)，即 \(\mathcal M\) 为因子。按照 \(\mathcal M\) 中投影的性质，定义下列类型。
\begin{center}
\begin{tabularx}{\linewidth}{@{}lX@{}}
\toprule
类型 & 投影条件 \\
\midrule
\(\mathrm I\) 型 & 存在非零极小投影。 \\
\(\mathrm{II}\) 型 & 不存在非零极小投影，但存在非零有限投影。 \\
\(\mathrm{III}\) 型 & 不存在非零有限投影。 \\
\bottomrule
\end{tabularx}
\end{center}
对于 \(\mathrm{II}\) 型因子，若单位投影有限，则称为 \(\mathrm{II}_1\) 型；若单位投影无限，则称为 \(\mathrm{II}_\infty\) 型。
\end{definition}

三种条件互相排斥并穷尽所有因子。它们区分的是投影结构，并未给出同一类型内部的同构分类。这里采用因子的投影判据；一般 von Neumann 代数的类型通过中心投影分解定义，各个非零中心分量可以属于不同类型。投影等价、有限性和一般的中心分解见 Blackadar\cite[\nopp III.1.1.2; Definition III.1.3.1; III.1.4.5--III.1.4.7]{Blackadar2017OperatorAlgebrasCorrected}。

\ProseParagraphBreak
由\cref{b2:prop:finite-projection-models}，\(B(H)\) 总是 \(\mathrm I\) 型因子。\(M_n(\mathbb C)\) 属于 \(\mathrm I_n\) 型，而 \(B(\ell^2(\mathbb N))\) 属于 \(\mathrm I_\infty\) 型；前者的单位投影有限，后者的单位投影无限。一般 \(\mathrm I\) 型因子也与某个 \(B(K)\) 同构，其结构证明是把等价的极小投影组织成矩阵单位，见\cite[\nopp III.1.5.2--III.1.5.3; III.1.5.12]{Blackadar2017OperatorAlgebrasCorrected}。当 \(K\) 不可分时，还可用其 Hilbert 维数作为 \(\mathrm I\) 型的下标。

\ProseParagraphBreak
\(\mathrm{II}_1\) 型也有有限单位投影，却没有矩阵代数中的极小投影；有限投影可以继续分解，不能把它理解成有限秩投影。\(\mathrm{II}_\infty\) 型的单位投影无限，但代数内部仍有非零有限投影；\(\mathrm{III}\) 型连这样的有限部分也不存在。尤其不能仅由“没有迹态”断定一个因子为 \(\mathrm{III}\) 型：\(B(\ell^2(\mathbb N))\) 就没有迹态。若它有迹态 \(\tau\)，取 \(v_1e_j=e_{2j}\)、\(v_2e_j=e_{2j-1}\)，则 \(v_i^*v_i=I\)，两个终投影之和为 \(I\)，因而 \(1=\tau(I)=\tau(v_1v_1^*)+\tau(v_2v_2^*)=2\)，矛盾。\(\mathrm{II}\) 与 \(\mathrm{III}\) 型的具体构造、非有界迹以及 \(\mathrm{III}\) 型的进一步细分见\cite[\nopp III.2.5; III.3; III.4.6]{Blackadar2017OperatorAlgebrasCorrected}。

% !TeX root = ../../Book2.tex
% 纲要标题：### D.5 与群、张量和非交换几何的联系

% 纲要要点（供目录核对）：
% * 从离散群代数 \(\mathbb C[G]\) 到全群／约化群 \(C^*\)-代数及群 von Neumann 代数；以 \(G=\mathbb Z\) 比较 Laurent 多项式与圆周函数
% * 代数张量积与算子代数张量积需要区分；定义最小／最大 \(C^*\)-张量范数，引用其极值定理并说明含幺情形的最大性
% * 高阶算子 K 理论、Tomita–Takesaki 理论、非交换几何和非自伴算子代数列为专题出口；本附录不是这些专著的压缩替代；作为本节末阅读说明，不另设小节
% #### D.5.1 离散群的算子代数
% #### D.5.2 代数张量积与算子张量范数
% ---

\section{与群、张量和非交换几何的联系}
\label{b2:sec:D05}

同一个群代数可以有不同的完备化，同一个代数张量积也可以有不同的算子范数。本节只建立辨认这些对象所需的定义，并在整数群及矩阵代数上计算。这样既能看到与本卷代数构造的联系，也能说明后续分析问题究竟增加了什么。

\subsection{离散群的算子代数}
\label{b2:subsec:D05:01}

对离散群 \(G\)，群代数 \(\mathbb C[G]\) 的元素是有限和 \(a=\sum_g a_g g\)。规定
\[
a^*=\sum_g\overline{a_g}g^{-1}.
\]
任何酉表示 \(U:G\to\mathcal U(H)\) 通过线性延拓给出 \(*\)-表示 \(U(a)=\sum_g a_g U_g\)，且
\[
\bigl\|U(a)\bigr\|\le\sum_g|a_g|.
\]
所以 \(\|a\|_{\max}=\sup_U\bigl\lVert U(a)\bigr\rVert\) 有限。它不只是一项半范数：在 \(\ell^2(G)\) 的标准正交基上定义左正则表示 \(\lambda_g\delta_h=\delta_{gh}\)，则
\(\lambda(a)\delta_e=\sum_g a_g\delta_g\)，故 \(a\ne0\) 时 \(\lambda(a)\ne0\)。次乘性、对合及 \(C^*\) 恒等式由各算子表示继承，得到一个真正的 \(C^*\) 范数。

\begin{definition}\label{b2:def:group-operator-algebras}
设 \(G\) 为离散群，\(\mathbb C[G]\) 为其复群代数。对 \(a=\sum_g a_g g\in\mathbb C[G]\)，令 \(\|a\|_{\max}=\sup_U\|\sum_g a_gU_g\|\)，其中 \(U\) 遍历 \(G\) 在复 Hilbert 空间上的酉表示。将 \(\mathbb C[G]\) 按此范数完备化，所得代数称为\term[quan-qun-Cxing-daishu]{全群 \(C^*\)-代数}[full group C*-algebra]，记为 \(C^*(G)\)。在 \(\ell^2(G)\) 的标准基上令 \(\lambda_g\delta_h=\delta_{gh}\)。将左正则表示 \(\lambda:\mathbb C[G]\to B(\ell^2(G))\) 的像取算子范数闭包，所得代数称为\term[yuehua-qun-Cxing-daishu]{约化群 \(C^*\)-代数}[reduced group C*-algebra]，记为 \(C_r^*(G)\)；若取弱算子拓扑下的闭包，则称所得代数为\term[qun-von-Neumann-daishu]{群 von Neumann 代数}[group von Neumann algebra]，记为 \(L(G)\)。
\end{definition}
\symidx{\(C^*(G)\)：离散群的全群 \(C^*\)-代数}
\symidx{\(C_r^*(G)\)：离散群的约化群 \(C^*\)-代数}
\symidx{\(L(G)\)：群 von Neumann 代数}

\(L(G)=\lambda(G)''\) 由\cref{b2:thm:operator-bicommutant}成立。每个酉表示唯一延拓到 \(C^*(G)\)，因为它对最大范数收缩；但不能不加条件地断言它对约化范数也收缩。全群与约化群代数间的自然映射及其差异见 Blackadar\cite[\nopp II.10.2.3--II.10.2.7]{Blackadar2017OperatorAlgebrasCorrected}。

\begin{proposition}\label{b2:prop:integer-group-cstar}
设 \(G=\mathbb Z\)，以 \(u\) 表示群代数 \(\mathbb C[G]\) 中对应于 \(1\in G\) 的元素，并令 \(u^*=u^{-1}\)。记单位圆为 \(\mathbb T=\bigl\{\,z\in\mathbb C\,\bigm|\,|z|=1\,\bigr\}\)，\(C(\mathbb T)\) 为其上的复连续函数代数，则 \(\mathbb C[G]=\mathbb C[u,u^{-1}]\)，且
\[
C^*(\mathbb Z)\cong C_r^*(\mathbb Z)\cong C(\mathbb T).
\]
\end{proposition}
\begin{proof}
整数群的酉表示由一个酉算子 \(U\) 决定。其谱包含于 \(\mathbb T\)：\(|z|>1\) 时用 Neumann 级数，\(|z|<1\) 时写 \(U-zI=U(I-zU^*)\)，仍可求逆。故对 Laurent 多项式 \(p\)，\cref{b2:thm:continuous-functional-calculus}给出
\(\bigl\lVert p(U)\bigr\rVert\le\max_{|z|=1}\bigl|p(z)\bigr|\)。反过来，每个标量 \(z\in\mathbb T\) 都给出一维酉表示，所以最大范数正好为这个上确界。

圆周弧长测度 \(\mu\) 取总质量一。函数 \(z^n\)、\(n\in\mathbb Z\)，两两正交且各有范数一，因为 \(\int_{\mathbb T}z^j\,d\mu(z)=0\) 对非零整数 \(j\) 成立。Laurent 多项式由 Stone--Weierstrass 定理在连续函数中一致稠密，连续函数又在 \(L^2\) 中稠密，所以这些函数构成 \(L^2(\mathbb T)\) 的正交基。因此 \(F\delta_n=z^n\) 从有限支集向量的等距映射延伸为酉同构 \(F:\ell^2(\mathbb Z)\to L^2(\mathbb T)\)。这证明所需的\term[Fourier-Plancherel-dingli]{Fourier--Plancherel 定理}[Fourier--Plancherel theorem]的圆周版本，亦见\cite[Theorem II.10.1.15(iii)]{Blackadar2017OperatorAlgebrasCorrected}。

由 \(F\lambda_1\delta_n=z^{n+1}\)，双边移位在此同构下恰为乘坐标函数。乘连续函数 \(p\) 的算子范数为 \(\max|p|\)：上界直接成立，下界可取最大值附近正测度弧段上的归一化示性函数，使像的范数任意接近最大值。因此约化范数与最大范数一致。最后，Laurent 多项式的一致稠密性使完备化得到 \(C(\mathbb T)\)。
\end{proof}

同一个 Fourier 模型还给出 \(L(\mathbb Z)\cong L^\infty(\mathbb T)\)。证明可由交换子完成。若有界算子 \(T\) 与乘坐标函数及其逆交换，它就与全部 Laurent 多项式、继而与全部连续函数的乘法算子交换。弧长测度的正则性使每个可测集合 \(E\) 的示性函数都可由连续函数 \(f_n\) 在 \(L^2\) 中逼近，且可取 \(0\le f_n\le1\)。记乘函数 \(h\) 的算子为 \(M_h\)。对 \(\xi\in L^\infty(\mathbb T)\subset L^2(\mathbb T)\)，有
\[
\|(M_{f_n}-M_{1_E})\xi\|_2
\le\|\xi\|_\infty\|f_n-1_E\|_2\longrightarrow0.
\]
对一般 \(\xi\in L^2(\mathbb T)\)，先用截断将它逼近为有界函数 \(\xi_k\)。由于 \(\|M_{f_n}-M_{1_E}\|\le1\)，
\[
\|(M_{f_n}-M_{1_E})\xi\|_2
\le\|\xi-\xi_k\|_2
+\|(M_{f_n}-M_{1_E})\xi_k\|_2.
\]
先取 \(k\) 再取 \(n\)，即得 \(M_{f_n}\to M_{1_E}\) 的强算子收敛。于是 \(T\) 也与每个示性函数的乘法算子交换。令 \(g=T1\)，对可测集合 \(E\) 有 \(T1_E=1_Eg\)。于是
\[
\|1_Eg\|_2\le\|T\|\cdot\|1_E\|_2.
\]
若 \(|g|>\|T\|\) 的集合有正测度，取其中 \(|g|\ge\|T\|+\varepsilon\) 的正测度子集，便违反此式。故 \(g\in L^\infty\)，\(T\) 在简单函数及其 \(L^2\) 极限上都是乘 \(g\)。同样的论证表明，全部本质有界函数的乘法代数也是自身的交换子。因此乘连续函数代数的交换子和双交换子均为 \(L^\infty\) 乘法代数；\cref{b2:thm:operator-bicommutant}和 Fourier 同构便给出结论。它的元素按几乎处处相等认同，包含弧段的示性函数；这些通常不是连续函数。于是代数有限和、范数完备化与弱算子闭包确实是三个不同层次。一般交换群的 Fourier 变换与群 \(C^*\)-代数分别见\cite[Theorem II.10.1.15; Proposition II.10.2.7]{Blackadar2017OperatorAlgebrasCorrected}；\(L(G)\) 与 \(L^\infty(\widehat G)\) 的联系及其 Hopf--von Neumann 背景另见同书 II.10.8.13。

\ProseParagraphBreak
一个非交换的例子来自两个生成元上的自由群 \(F_2\)：\(L(F_2)\) 是 \(\mathrm{II}_1\) 因子，具有忠实迹态 \(\tau(T)=\langle\delta_e,T\delta_e\rangle\)，却没有非零极小投影。其因子性与 \(F_2\) 中每个非单位元的共轭类均无限有关：中心元的平方可和系数在共轭类上恒定，因而只能支撑于单位元。迹与投影类型的完整论证见\cite[\nopp III.3.3.2--III.3.3.7]{Blackadar2017OperatorAlgebrasCorrected}。

\subsection{代数张量积与算子张量范数}
\label{b2:subsec:D05:02}

第六章的张量积只有有限和。本小节固定两个非零复 \(C^*\)-代数 \(A,B\)，用 \(A\odot B\) 特别表示它们在复数域上的代数张量积，规定 \((a\otimes b)^*=a^*\otimes b^*\)。符号改变只是为了与下面的完备化区分。
\symidx{\(A\odot B\)：两个复代数的代数张量积}

\ProseParagraphBreak
取 \(A,B\) 的忠实 \(*\)-表示 \(\pi_A:A\to B(H_A)\)、\(\pi_B:B\to B(H_B)\)，先构造表示所作用的空间。在复向量空间的代数张量积 \(H_A\odot H_B\) 上规定
\[
\langle\xi\otimes\eta,\xi'\otimes\eta'\rangle
=\langle\xi,\xi'\rangle\langle\eta,\eta'\rangle,
\]
并对第一变量共轭线性、对第二变量线性地延拓。任意有限张量都落在两侧某个有限维子空间的张量积中；取这两个子空间的正交归一基后，张量积基也正交归一，所得平方范数就是系数绝对值的平方和，所以这个内积正定。按内积范数完备化，得到 Hilbert 张量积 \(H_A\otimes H_B\)。有界算子 \(S\in B(H_A)\)、\(T\in B(H_B)\) 在有限张量上按 \(\xi\otimes\eta\mapsto S\xi\otimes T\eta\) 作用；将有限张量按一侧正交归一基展开，分别作用于两侧，得到范数上界 \(\|S\|\cdot\|T\|\)，故作用连续延伸到这个 Hilbert 张量积。

\ProseParagraphBreak
现在换到代数 \(A\odot B\) 上：令 \(a\otimes b\) 在已经完成的 Hilbert 空间 \(H_A\otimes H_B\) 上按 \(\pi_A(a)\otimes\pi_B(b)\) 作用，并对有限和线性延拓。所得算子范数称为\term[kongjian-zhangliang-fanshu]{空间张量范数}[spatial tensor norm]，亦记 \(\|\cdot\|_{\min}\)。随后按这个范数完备化的是 \(A\odot B\)，得到 \(A\otimes_{\min}B\)。前一次完备化得到表示空间，后一次得到 \(C^*\)-代数；它们作用于不同的对象。

\ProseParagraphBreak
另一方面，另取复 Hilbert 空间 \(H\)，在其上取像互相交换的 \(*\)-表示 \(\pi:A\to B(H)\)、\(\rho:B\to B(H)\)，考察
\[
\sum_i a_i\otimes b_i\longmapsto\sum_i\pi(a_i)\rho(b_i).
\]
记这个诱导的 \(*\)-表示为 \(\pi\cdot\rho\)。对所有这种表示取范数上确界，得到\term[zuidazhangliang-fanshu]{最大张量范数}[maximal tensor norm] \(\|\cdot\|_{\max}\)。\cref{b2:prop:cstar-star-hom-general}保证两个表示均收缩，因而
\[
\left\|\sum_i\pi(a_i)\rho(b_i)\right\|
\le\sum_i\|a_i\|\,\|b_i\|,
\]
所取上确界有限。按这个范数完备化 \(A\odot B\)，所得代数记为 \(A\otimes_{\max}B\)。最小、最大范数比较的对象都是代数 \(A\odot B\)，属于后一次完备化的选择。
\symidx{\(A\otimes_{\min}B\)：空间张量积的范数完备化}
\symidx{\(A\otimes_{\max}B\)：最大张量积的范数完备化}

\ProseParagraphBreak
忠实表示的存在性由\cref{b2:thm:gelfand-naimark-unital,b2:prop:cstar-unitization-representation}保证。矩阵系数能分离有限张量，故 \(\pi_A\otimes\pi_B\) 在 \(A\odot B\) 上单射，空间范数确为范数。两个算子因子的作用彼此交换，所以空间表示也是最大范数定义中容许的一种表示；这同时保证最大范数的非退化性。次乘性与对合的等距性由各表示继承，而
\[
\|x^*x\|_{\max}
=\sup_{\pi,\rho}\|(\pi\cdot\rho)(x)\|^2
=\|x\|_{\max}^2
\]
给出 \(C^*\) 恒等式。

\begin{claim}[最小与最大 \(C^*\) 张量范数]\label{b2:claim:cstar-tensor-extrema}
设 \(A,B\) 为非零复 \(C^*\)-代数，\(A\odot B\) 为其复代数张量积。取忠实表示 \(\pi_A:A\to B(H_A)\)、\(\pi_B:B\to B(H_B)\)，以 \(\pi_A\otimes\pi_B\) 在 \(H_A\otimes H_B\) 上的算子范数定义 \(\|\cdot\|_{\min}\)；对同一 Hilbert 空间上像互相交换的全部 \(*\)-表示 \(\pi,\rho\)，以 \(x=\sum_i a_i\otimes b_i\) 对应算子 \(\sum_i\pi(a_i)\rho(b_i)\) 的范数上确界定义 \(\|x\|_{\max}\)。空间范数不依赖所选的两个忠实表示。对 \(A\odot B\) 上任意 \(C^*\) 范数 \(\gamma\)，都有
\[
\|x\|_{\min}\le\gamma(x)\le\|x\|_{\max}
\qquad(x\in A\odot B),
\]
且 \(\gamma(a\otimes b)=\|a\|\,\|b\|\) 对所有 \(a\in A,b\in B\) 成立。
\end{claim}

这项一般定理见 Blackadar\cite[\nopp II.9.1.3; II.9.5.1--II.9.5.2]{Blackadar2017OperatorAlgebrasCorrected}。最大性的含幺情形可以直接说明：设 \(C=A\otimes_\gamma B\)，在 \(C\) 的忠实 Hilbert 空间表示中，映射 \(a\mapsto a\otimes1_B\)、\(b\mapsto1_A\otimes b\) 是像互相交换的 \(*\)-同态，由\cref{b2:prop:cstar-star-hom-general}自动收缩。因此 \(C\) 的算子范数受最大范数控制，即 \(\gamma\le\|\cdot\|_{\max}\)。

极小性所需的分析工具更多。一个态若不能写成两个不同态的非平凡凸组合，就称为纯态。标准证明先建立纯态与不可约 GNS 表示的对应，再证明纯态乘积能延拓到 \(A\otimes_\gamma B\) 上的态，由这些表示控制空间范数。这里用到态空间的弱星紧性、紧凸集的极点与分离定理，以及交换因子张量积的理想结构；非含幺情形还要用近似单位。相关结果分别见同书 II.6.3--II.6.4、Theorem II.9.4.4 和 II.9.5.1；非含幺表示的处理见 Theorem II.9.2.1 与 Corollary II.9.2.2。这些输入超出此前的 GNS 构造，代数张量积的双线性泛性质本身并不能推出上述范数比较。

在矩阵例子中，\(E_{ij}\otimes E_{ab}\) 送到 \(\mathbb C^n\otimes\mathbb C^m\) 上相应的矩阵单位，给出
\[
M_n(\mathbb C)\odot M_m(\mathbb C)\cong M_{nm}(\mathbb C).
\]
这一代数同构也保持共轭转置；空间范数就是右边的算子范数。最大范数也相同：任何 \(C^*\) 范数 \(\gamma\) 都使这个有限维空间完备，由\cref{b2:thm:continuous-functional-calculus}作用于正矩阵 \(x^*x\)，其范数为最大特征值，所以 \(\gamma(x)^2=\gamma(x^*x)\) 就是通常算子范数的平方。有限维空间已经完备，因此该例不增加新元素。相比之下，\(C\bigl([0,1]\bigr)\odot C\bigl([0,1]\bigr)\) 的元素仅是有限和 \(\sum_i f_i(x)g_i(y)\)；这些函数含常数、分离正方形上的点且对共轭封闭，故由 Stone--Weierstrass 定理在一致范数下稠密于 \(C\bigl([0,1]^2\bigr)\)，但不包含每一个二元连续函数。\cref{b2:ex:D-tensor-completion}给出这样的具体函数及一致逼近，说明完备化在这里有实际作用。

\ProseParagraphBreak

Gelfand 对偶允许把交换函数代数当作拓扑空间的一种代数描述。将其推广为“非交换空间”的研究方法时，需要先指定保留哪些数据：\(C^*\)-代数重视范数与连续表示，von Neumann 代数重视弱闭包与测度结构，附录C的 Azumaya 理论则固定交换基环并研究局部的代数矩阵结构。这些问题有联系，却不能仅凭名称就认作同一种几何。

\ProseParagraphBreak
算子 K 理论通过稳定矩阵中的投影、可逆元及连续同伦研究代数；第四卷附录B的代数 K 理论初步只能提供部分比较语言。Tomita--Takesaki 理论\footnote{\OriginalHanName{冨田稔}（假名：とみた みのる，罗马音：Minoru Tomita，1924-2015），日本数学家，建立 von Neumann 代数的模理论。竹崎正道（假名：たけさき まさみち，罗马音：Masamichi Takesaki，1933-），日本数学家，系统阐述并发展这一理论。}从 von Neumann 代数的态和权构造模自同构群；这里的“模”是该理论的专名，不是本卷任意环模的简称。


\begin{chaptersummary}
\(C^*\) 恒等式把代数平方与算子大小联系起来，谱论和连续函数演算给出 \(A_+=\{x^*x:x\in A\}\)。Hahn--Banach 延拓先保留范数和单位元处的取值，再由估计推出正性。GNS 用半内积的零长度左理想构造向量空间商，左乘稳定性保证作用良定义，同一个估计再保证作用有界；Hilbert 完备化产生表示。态的忠实性与表示的忠实性不同，矩阵代数的非忠实向量态仍能给出忠实的自然表示；全部态的 Hilbert 直和则检测任意抽象代数元素。

\ProseParagraphBreak
双交换子定理说明，含幺自伴算子代数的全部交换关系决定其强、弱算子闭包。证明必须同时处理有限多个向量，不能只给每个向量分别选择一个逼近算子。范数闭包通常更小，故 \(C^*\)-代数与 von Neumann 代数保留的信息不同。因子的类型进一步取决于代数内部投影的等价关系：有无极小投影、有无非零有限投影，将 I、II、III 型区分开来；有限性并不等于所表示子空间的有限维性。

\ProseParagraphBreak
整数群把这种差异具体化为有限 Laurent 和、\(C(\mathbb T)\) 与 \(L^\infty(\mathbb T)\)。一般离散群的全群代数自然满射到约化群代数，两者使用的表示范围不同。张量构造则有两次不同的完备化：先完成表示空间的内积范数得到 Hilbert 张量积，再按最小或最大 \(C^*\) 范数完成代数的有限张量和。矩阵张量积本已有限维，不增加元素；函数张量积却可增加无法写成有限分离和的连续函数。最大范数使每对像交换的有界 \(*\)-表示唯一延拓，最小性及表示无关性仍需所列的分析输入。附录C的 Azumaya 局部矩阵结构固定交换基环，这里不同闭包所保留的拓扑与测度信息须另行确定。
\end{chaptersummary}

\BookExercises

\begin{exercise}\label{b2:ex:D-norm-and-involution}
在 \(\mathbb C\) 上取范数 \(\|z\|'=2|z|\)，证明它满足 Banach 代数的条件，却不满足通常复共轭下的 \(C^*\) 恒等式。再说明 \(M_2(\mathbb C)\) 的普通转置为何不是本附录定义的共轭线性对合。
\end{exercise}
\begin{solution}
完备性不因正比例缩放改变，且 \(2|zw|\le\bigl(2|z|\bigr)\cdot\bigl(2|w|\bigr)\)。但对 \(z=1\)，\(\|\bar z z\|'=2\ne4=\|z\|'^2\)。普通转置满足 \((iI)^{\mathsf T}=iI\)，共轭线性则要求它为 \(-iI\)，故不符合该条件。它仍是复代数上的第一类线性对合，见\cref{b2:def:B-algebra-involution}。
\end{solution}

\begin{exercise}\label{b2:ex:D-diagonal-compact}
设 \((w_j)\) 为有界复数列，在 \(H=\ell^2(\mathbb N)\) 的标准正交基上令 \(We_j=w_je_{j+1}\)。证明 \(\|W\|=\sup_j|w_j|\)，并求只保留前 \(n\) 个权重的截断与 \(W\) 的范数距离。证明 \(W\) 紧当且仅当 \(w_j\to0\)，且
\[
\inf_{K\in\mathcal K(H)}\|W-K\|=\limsup_{j\to\infty}|w_j|.
\]
\end{exercise}
\begin{solution}
由 \(\|W\xi\|^2=\sum_j|w_j|^2|\xi_j|^2\)，逐项估计得范数上界；在 \(e_j\) 处取值给出下界，所以范数为权重绝对值的上确界。若 \(W_n\) 保留前 \(n\) 个权重，则同样计算得到 \(\|W-W_n\|=\sup_{j>n}|w_j|\)。因而权重趋零时，有限秩截断在范数下收敛于 \(W\)，使 \(W\) 紧。

对任意紧算子 \(K\)，有 \(Ke_j\to0\)。确实，有限秩算子可以写成有限和 \(F\xi=\sum_i\langle v_i,\xi\rangle u_i\)，每个 \(\langle v_i,e_j\rangle\) 都趋零；再用有限秩算子在范数下逼近 \(K\)，即可使 \(\|Ke_j\|\) 趋零。因此
\[
\|W-K\|\ge\|We_j-Ke_j\|\ge|w_j|-\|Ke_j\|,
\]
取上极限给出所求距离的下界。另一方面，取 \(K=W_n\)，令 \(n\to\infty\)，尾部上确界趋于 \(\limsup_j|w_j|\)，给出反向不等式。若 \(W\) 本身紧，这个距离为零，遂有 \(w_j\to0\)。权重有无限多个非零项时，其非零像向量线性无关，所以紧算子仍可有无限秩。
\end{solution}

\begin{exercise}\label{b2:ex:D-projection-versus-idempotent}
设 \(u,v\in\mathbb C^n\) 满足 \(v^*u=1\)，令 \(E=uv^*\)，并令 \(Q\) 为到 \(\mathbb Cu\) 的正交投影。证明 \(E^2=E\)，求 \(\|E\|\) 和 \(\|E-Q\|\)，并证明 \(E\) 自伴当且仅当 \(\|E\|=1\)。
\end{exercise}
\begin{solution}
\(E^2=u(v^*u)v^*=E\)。由 \(E\xi=u\langle v,\xi\rangle\) 及 Cauchy--Schwarz，\(\|E\|\le\|u\|\,\|v\|\)，在 \(\xi=v/\|v\|\) 处达到，所以等号成立。正交投影为 \(Q=uu^*/\|u\|^2\)。写 \(w=v-u/\|u\|^2\)，由 \(u^*v=1\) 得 \(u^*w=0\)。于是
\[
E-Q=uw^*,\qquad
\|E-Q\|^2=\|u\|^2\|w\|^2=\|u\|^2\|v\|^2-1.
\]
因此 \(\|E\|=1\) 当且仅当 \(w=0\)，也就是 \(E=Q\)。若 \(E\) 自伴，它是到自身像 \(\mathbb Cu\) 的正交投影，故必为 \(Q\)；反向显然。这说明幂等矩阵的范数超过一，具体测量了它与同像正交投影的偏差。
\end{solution}

\begin{exercise}\label{b2:ex:D-finite-function-maps}
分类全部保持单位元的 \(*\)-同态 \(\mathbb C^r\to\mathbb C^s\)，并判断它们何时单射、何时满射。
\end{exercise}
\begin{solution}
将同态与目标的第 \(j\) 个坐标取值复合，得到源的一个特征标 \(\chi_j\)。记源的坐标幂等元为 \(e_i\)：\(\chi_j(e_i)^2=\chi_j(e_i)\) 使各值为零或一，\(e_ie_\ell=0\)、\(i\ne\ell\)，使至多一个值为一，而 \(\sum_i e_i=1\) 使恰有一个值为一。把这个指标记为 \(f(j)\)，则 \(\chi_j(a)=a_{f(j)}\)。所以原同态为 \((a_i)\mapsto\bigl(a_{f(j)}\bigr)\)，由函数 \(f:\{1,\ldots,s\}\to\{1,\ldots,r\}\) 唯一决定。每个源坐标都出现时且仅此时映射单射，即 \(f\) 满射；各目标坐标可独立指定时且仅此时映射满射，即 \(f\) 单射。反向的公式保持乘法、单位和逐点共轭。
\end{solution}

\begin{exercise}\label{b2:ex:D-czero-pullback}
证明 \(f(x)=x^2\) 的拉回把 \(C_0(\mathbb R)\) 映入自身，而常值映射 \(g(x)=0\) 的拉回没有这个性质。
\end{exercise}
\begin{solution}
紧集在平方映射下的原像闭且有界，故紧。对 \(h\in C_0(\mathbb R)\) 及 \(\varepsilon>0\)，满足 \(|h\circ f|\ge\varepsilon\) 的点，组成 \(h\) 的相应紧超水平集在 \(f\) 下的原像，所以拉回仍趋零。取 \(h(t)=e^{-t^2}\)，则 \(h\circ g=1\)，不属于 \(C_0(\mathbb R)\)。两个映射都连续，差别在紧集原像的性质。
\end{solution}

\begin{exercise}\label{b2:ex:D-functional-calculus-matrix}
令 \(A=\begin{psmallmatrix}2&1\\1&2\end{psmallmatrix}\)。求谱及两个谱投影，写出任意谱上函数 \(f\) 的 \(f(A)\)，并求正平方根。
\end{exercise}
\begin{solution}
单位特征向量为 \((1,1)^{\mathsf T}/\sqrt2\)、\((1,-1)^{\mathsf T}/\sqrt2\)，特征值分别为三、一。相应投影为
\[
P_3=\frac12\begin{pmatrix}1&1\\1&1\end{pmatrix},\qquad
P_1=\frac12\begin{pmatrix}1&-1\\-1&1\end{pmatrix}.
\]
所以 \(f(A)=f(3)P_3+f(1)P_1\)，正平方根为 \(\sqrt3P_3+P_1\)。投影正交且和为单位矩阵，直接平方即可核验；正性来自两个非负特征值。
\end{solution}

\begin{exercise}\label{b2:ex:D-density-matrices}
令 \(h=\operatorname{diag}(4,4,-1)\in M_3(\mathbb C)\)。求态 \(\varphi\) 上 \(\varphi(h)\) 的最大值，描述全部达到最大值的态。进一步证明，在这组态中，不能写成两个不同态的平均的元素恰为相应秩一密度矩阵给出的态。
\end{exercise}
\begin{solution}
由\cref{b2:prop:D-density-gns}，态唯一写成 \(\varphi(a)=\operatorname{tr}(\rho a)\)，其中 \(\rho\ge0\)、\(\operatorname{tr}\rho=1\)。故 \(\varphi(h)=4-5\rho_{33}\le4\)，在 \(\rho_{33}=0\) 时且仅此时达到四。半正定矩阵的 \(2\times2\) 主子式给出 \(|\rho_{i3}|^2\le\rho_{ii}\rho_{33}=0\)，所以最后一行、列都为零。达到最大值的密度矩阵恰为
\[
\rho=\begin{pmatrix}D&0\\0&0\end{pmatrix},
\qquad D\in M_2(\mathbb C),\quad D\ge0,\quad\operatorname{tr}D=1.
\]
若 \(D\) 秩二，取两个正特征值 \(\lambda_1,\lambda_2\)，选 \(0<\varepsilon<\min(\lambda_1,\lambda_2)\)，在酉特征基中把它们分别改为 \((\lambda_1\pm\varepsilon,\lambda_2\mp\varepsilon)\)。得到两个不同的半正定迹一矩阵，原矩阵是其平均，所以相应态不是极点。

若 \(\rho=vv^*\)、\(\|v\|=1\) 且 \(v\) 在前两个坐标中，假设它为两个密度矩阵 \(\rho_1,\rho_2\) 的平均。对任意 \(w\perp v\)，有 \(0=w^*\rho w=(w^*\rho_1w+w^*\rho_2w)/2\)，两个非负项均为零。由正平方根，\(w^*\rho_iw=\|\rho_i^{1/2}w\|^2=0\)，故 \(\rho_iw=0\)。两矩阵的像都在 \(\mathbb Cv\)，迹一便使它们都等于 \(vv^*\)。因此这样的态不能由两个不同态平均得到；论证还排除了在全部态中作这样的分解。
\end{solution}

\begin{exercise}\label{b2:ex:D-density-gns}
在 \(A=M_2(\mathbb C)\oplus\mathbb C\) 上取态
\[
\varphi_t(a,z)=t\langle e_1,ae_1\rangle+(1-t)z,
\qquad0\le t\le1.
\]
显式写出它的 GNS 空间、作用和循环向量，求表示的核与维数。判断态本身和表示各在何时忠实。
\end{exercise}
\begin{solution}
对 \((a,z)\in A\)，有 \(\varphi_t((a,z)^*(a,z))=t\|ae_1\|^2+(1-t)|z|^2\)。若 \(0<t<1\)，映射
\[
[(a,z)]\longmapsto(\sqrt t\,ae_1,\sqrt{1-t}\,z)
\]
将 GNS 空间酉识别为 \(\mathbb C^2\oplus\mathbb C\)，左作用为 \((b,w)(\xi,c)=(b\xi,wc)\)，循环向量为 \((\sqrt t\,e_1,\sqrt{1-t})\)。像取遍该三维空间，且每个非零分量都被相应作用检测，所以表示核为零。

在 \(t=0\) 时，空间为 \(\mathbb C\)，作用只取标量分量，循环向量为一，核为 \(M_2(\mathbb C)\oplus0\)。在 \(t=1\) 时，空间为 \(\mathbb C^2\)，作用只取矩阵分量，循环向量为 \(e_1\)，核为 \(0\oplus\mathbb C\)。对应维数分别为一、二。因而表示忠实当且仅当 \(0<t<1\)。态本身在全部参数下都不忠实，因为非零元素 \((E_{22},0)\) 的正平方取值恒为零；中间参数因此给出态不忠实而表示忠实的例子。
\end{solution}

\begin{exercise}\label{b2:ex:D-finite-gns}
在 \(A=\mathbb C^r\) 上令 \(\varphi(a)=\sum_j t_j a_j\)，其中 \(t_j\ge0\)、\(\sum_jt_j=1\)。显式写出 GNS 空间、作用和循环向量，判断态和表示各在何时忠实。
\end{exercise}
\begin{solution}
零长度子空间由在所有 \(t_j>0\) 处取零的向量组成。记这些指标为 \(J\)，则商为 \(\mathbb C^J\)，内积是 \(\sum_{j\in J}t_j\bar a_j b_j\)，作用为逐点乘法，循环向量为全一向量。若用标准内积，则等距识别将循环向量改为 \((\sqrt{t_j})_{j\in J}\)，作用仍对角。表示忠实当且仅当没有丢失任何坐标，即全部权重都正。由 \(\varphi(a^*a)=\sum_jt_j|a_j|^2\)，态忠实也恰有同一条件；这个交换例子中二者一致，矩阵代数的向量态则未必如此。
\end{solution}

\begin{exercise}\label{b2:ex:D-positive-not-hom}
在 \(M_2(\mathbb C)\) 上令 \(\mathcal E(a)=(a+a^{\mathsf T})/2\)。证明 \(\mathcal E\) 复线性、保单位、保正且 \(\mathcal E^2=\mathcal E\)。求它的像与核，并判断它的像是否为子代数、它是否为代数同态。
\end{exercise}
\begin{solution}
转置保正由\cref{b2:prop:D-positive-not-hom}给出；恒等映射也保正，正锥的凸性遂保证其平均保正。复线性、保单位和 \(\mathcal E^2=\mathcal E\) 直接由两次转置为恒等得到。像为对称矩阵空间，核为斜对称矩阵空间；在二阶情形，核是 \(\mathbb C(E_{12}-E_{21})\)。

像中有 \(E_{11}\) 和 \(E_{12}+E_{21}\)，二者乘积为不对称的 \(E_{12}\)，故像不是子代数。又
\[
\mathcal E(E_{12}E_{21})=E_{11},\qquad
\mathcal E(E_{12})\mathcal E(E_{21})=\tfrac14 I_2,
\]
所以 \(\mathcal E\) 不保持乘法。线性投影可以保正，却不一定把代数映到一个子代数。
\end{solution}

\begin{exercise}\label{b2:ex:D-multiplicity-commutant}
在 \(H=\mathbb C^n\otimes\mathbb C^m\) 上求 \(\mathcal A=M_n(\mathbb C)\otimes I_m\) 的交换子代数、双交换子和中心。
\end{exercise}
\begin{solution}
把 \(T\) 写成 \(n\times n\) 个 \(m\) 阶块。与 \(E_{ii}\otimes I_m\) 交换使非对角块全为零；再与 \(E_{ij}\otimes I_m\) 交换，使所有对角块相同。因此 \(\mathcal A'=I_n\otimes M_m(\mathbb C)\)。交换两因子的角色得到 \(\mathcal A''=\mathcal A\)，而两者的交只有标量，故中心为 \(\mathbb CI\)。重数 \(m\) 出现在交换子代数中，不改变 \(\mathcal A\) 自身的中心。
\end{solution}

\begin{exercise}\label{b2:ex:D-adjoint-strong}
设 \((U_\alpha)\) 为复 Hilbert 空间上的酉算子网，且它在 SOT 下收敛于酉算子 \(U\)。证明 \(U_\alpha^*\to U^*\) 也在 SOT 下成立。再在 \(\ell^2(\mathbb N)\) 上令 \(U_n\) 将前 \(n+1\) 个标准基向量循环置换：\(U_ne_j=e_{j+1}\) 对 \(j\le n\)，\(U_ne_{n+1}=e_1\)，其余基向量不动。求 \(U_n\) 的强极限，并判断伴随是否强收敛。
\end{exercise}
\begin{solution}
对每个 \(\xi\)，酉性给出
\[
\|(U_\alpha^*-U^*)\xi\|
=\|\xi-U_\alpha U^*\xi\|\longrightarrow0,
\]
因为 \(UU^*\xi=\xi\)。这里极限算子也酉，是把右边写成一个强收敛误差的条件。

每个 \(U_n\) 都把一组正交基置换成自身，故酉。对固定基向量，\(U_ne_j\) 最终为 \(e_{j+1}\)，因此在有限支集向量上收敛于单边移位 \(S\)；\(\|U_n\|=\|S\|=1\) 再把收敛延伸到全部向量，得到 \(U_n\to S\) 强收敛。可是 \(U_n^*e_1=e_{n+1}\)，这些向量彼此距离为 \(\sqrt2\)，没有范数极限，所以伴随网没有强极限。单边移位是等距但不满的算子，说明酉算子集合在 SOT 下不闭；前一断言的极限酉假设不能省略。
\end{solution}

\begin{exercise}\label{b2:ex:D-selfadjoint-hypothesis}
令 \(\mathcal A=\{\operatorname{diag}(z,0):z\in\mathbb C\}\subseteq B(\mathbb C^2)\)。求其交换子和双交换子，判断它是否 WOT 闭、是否对伴随封闭，以及为什么双交换子不等于其 WOT 闭包。再把作用空间限制到 \(\mathbb Ce_1\)，比较所得代数的双交换子。
\end{exercise}
\begin{solution}
与 \(E_{11}\) 交换的矩阵恰为对角矩阵，所以 \(\mathcal A'\) 是全部对角矩阵，\(\mathcal A''\) 也为全部对角矩阵。有限维中 WOT 是条目的通常拓扑，\(\mathcal A\) 由非对角条目和右下条目为零的闭条件定义，故 WOT 闭。共轭转置仍属于 \(\mathcal A\)，但整个空间的单位矩阵不在其中，缺少的是\cref{b2:thm:operator-bicommutant}的含单位条件。

在 \(\mathbb Ce_1\) 上限制作用以后，得到 \(\mathbb C\) 的全部标量算子，它已经含该子空间的恒等映射，交换子与双交换子都等于自身。若仍在原空间而给 \(\mathcal A\) 添加 \(I_2\)，则得到 \(\mathcal A+\mathbb CI_2\)，正好为原来的双交换子。单位元取决于所使用的作用空间。
\end{solution}

\begin{exercise}\label{b2:ex:D-no-tracial-state}
在 \(H=\ell^2(\mathbb N)\) 上令 \(Se_j=e_{j+1}\)，记 \(\mathcal T=C^*(I,S)\subset B(H)\) 为由 \(I,S,S^*\) 生成的范数闭代数。证明每个迹态在 \(\mathcal K(H)\) 上都为零，并通过下式构造 \(\mathcal T\) 上的一个迹态：
\[
\tau(T)=\lim_{N\to\infty}\frac1N\sum_{j=1}^N\langle e_j,Te_j\rangle.
\]
需证明极限对全部 \(T\in\mathcal T\) 存在。
\end{exercise}
\begin{solution}
令 \(p=I-SS^*\)，这是到 \(\mathbb Ce_1\) 的投影。每个迹态 \(\sigma\) 都满足 \(\sigma(p)=\sigma(I)-\sigma(SS^*)=\sigma(I)-\sigma(S^*S)=0\)。矩阵单位 \(E_{ij}=S^{i-1}p(S^*)^{j-1}\) 属于 \(\mathcal T\)，迹性质使 \(\sigma(E_{ii})=\sigma(p)=0\)，正泛函的 Cauchy--Schwarz 进而使 \(\sigma(E_{ij})=0\)。为核对这些矩阵单位的有限和在 \(\mathcal K(H)\) 中范数稠密，令 \(P_n\) 投影到前 \(n\) 个坐标。若 \(F\xi=\sum_i\langle v_i,\xi\rangle u_i\) 有限秩，则 \(P_nu_i\to u_i\)、\(P_nv_i\to v_i\) 使 \(P_nFP_n\to F\) 范数收敛，前者为有限矩阵单位的和。再用紧算子的有限秩逼近，得到所需稠密性，也得到 \(\mathcal K(H)\subseteq\mathcal T\)。态连续，所以全部紧算子上的值都为零。

记题中的有限平均为 \(\tau_N\)。它是向量态的平均，故为态且范数为一。关系 \(S^*S=I\) 将每个字约化成 \(S^p(S^*)^q\)、\(p,q\ge0\)。其对角取值在 \(p\ne q\) 时全为零，在 \(p=q\) 时除前 \(p\) 项为零外均为一。因此每个有限字和的 \(\tau_N\) 都收敛。对一般 \(T\in\mathcal T\)，选字多项式 \(a\) 使 \(\|T-a\|<\varepsilon\)，则
\[
|\tau_N(T)-\tau_M(T)|\le2\varepsilon+|\tau_N(a)-\tau_M(a)|.
\]
这证明 \(\tau_N(T)\) 是 Cauchy 列。所得极限线性、保正且在单位上为一，所以是态。

在忽略有限秩算子的商中，\(SS^*=S^*S=I\)，所以由 \(S,S^*\) 生成的字多项式彼此交换；即任意两个字多项式的交换子有限秩。有限秩算子作用于 \(e_j\) 的结果趋零，再用范数逼近便知每个紧算子 \(K\) 都满足 \(Ke_j\to0\)，所以 \(\tau_N(K)\to0\)。故 \(\tau(ab)=\tau(ba)\) 对字多项式成立，再由范数稠密性与乘法连续性延伸到全部 \(a,b\in\mathcal T\)。这给出迹态，而\cref{b2:prop:D-no-tracial-state}说明它不可能成为整个 \(B(H)\) 上的迹态。
\end{solution}

\begin{exercise}\label{b2:ex:D-cyclic-fourier}
设 \(C_n=\langle g\rangle\)，取本原 \(n\) 次单位根 \(\zeta\)，即其乘法阶恰为 \(n\)。证明
\[
e_j=\frac1n\sum_{r=0}^{n-1}\zeta^{-jr}g^r\qquad(0\le j<n)
\]
是正交自伴幂等元，和为一；由此求 \(\mathbb C[C_n]\) 的全群 \(C^*\) 范数。
\end{exercise}
\begin{solution}
有限几何级数给出 \(\sum_r\zeta^{ar}=0\)，除非 \(n\mid a\)，此时和为 \(n\)。展开乘积即得 \(e_ie_j=\delta_{ij}e_j\)、\(\sum_j e_j=1\)。取对合并把求和指标换成 \(-r\)，得到 \(e_j^*=e_j\)。又 \(ge_j=\zeta^j e_j\)，故 \(p(g)=\sum_jp(\zeta^j)e_j\)。每个酉表示将 \(e_j\) 送到正交投影，给出范数上界 \(\max_j\bigl|p(\zeta^j)\bigr|\)；一维表示 \(g\mapsto\zeta^j\) 给下界。因此该最大值正是全群范数，且代数为 \(\mathbb C^n\)，无需再增加完备化元素。
\end{solution}

\begin{exercise}\label{b2:ex:D-integer-cosine}
在 \(C^*(\mathbb Z)\) 中令 \(a=(u+u^{-1})/2\)。求 \(a\) 的谱及范数，并判断 \(1-a\) 是否正、是否可逆。
\end{exercise}
\begin{solution}
由\cref{b2:prop:integer-group-cstar}，它对应 \(e^{i\theta}\mapsto\cos\theta\)。连续函数代数中，\(f-\lambda\) 可逆当且仅当它处处不为零，因为此时倒数仍连续。因此谱为值域 \([-1,1]\)，范数为一。\(1-a\) 逐点非负，故正；它在 \(1\in\mathbb T\) 处为零，所以不可逆。
\end{solution}

\begin{exercise}\label{b2:ex:D-direct-sum-traces}
求 \(M_2(\mathbb C)\oplus\mathbb C\) 的全部迹态，并判断哪些忠实。
\end{exercise}
\begin{solution}
迹条件使非对角矩阵单位上的值为零，因为 \(E_{ii}E_{ij}=E_{ij}\)、\(E_{ij}E_{ii}=0\) 对 \(i\ne j\) 成立；又 \(\tau(E_{12}E_{21})=\tau(E_{21}E_{12})\) 使两个对角值相等。设第一块单位元的值为 \(t\)，正性与总单位元归一化使 \(0\le t\le1\)，从而全部迹态为
\[
\tau(a,z)=t\,\operatorname{tr}(a)/2+(1-t)z.
\]
这些泛函确为迹态。忠实性等价于两块权重都正，即 \(0<t<1\)；端点各零化一个非零直和因子。
\end{solution}

\begin{exercise}\label{b2:ex:D-tensor-completion}
在 \([0,1]^2\) 上令 \(F(x,y)=(1-xy/2)^{-1}\)，并令 \(F_N(x,y)=\sum_{n=0}^N2^{-n}x^ny^n\)。证明 \(F\) 不能写成有限和 \(\sum_{i=1}^r f_i(x)g_i(y)\)，但 \(\|F-F_N\|_\infty\le2^{-N}\)。证明 \(F_N\) 至少需要 \(N+1\) 个这样的分离项。
\end{exercise}
\begin{solution}
若 \(F\) 有 \(r\) 项分离表示，取 \(r+1\) 个不同的 \(y_j\in(0,1]\)。函数 \((1-y_jx/2)^{-1}\) 全在线性空间 \(\operatorname{span}\{f_1,\ldots,f_r\}\) 中，因此线性相关。假设系数为 \(c_j\)，乘去全部分母，得到在 \([0,1]\) 上恒为零的多项式
\[
\sum_j c_j\prod_{i\ne j}(1-y_ix/2)=0.
\]
它是零多项式，因而可以在区间外的 \(x=2/y_j\) 处取值，得到 \(c_j\prod_{i\ne j}(1-y_i/y_j)=0\)。因各 \(y_j\) 不同，乘积非零，全部 \(c_j=0\)，矛盾。

几何级数的余项为 \((xy/2)^{N+1}/(1-xy/2)\)，在正方形上不超过 \(2^{-N}\)，故所述一致逼近成立。对 \(F_N\)，取 \(N+1\) 个不同的 \(y_j\in[0,1]\)。这些切片在基 \(1,x,\ldots,x^N\) 下的系数矩阵是将 Vandermonde 矩阵各行乘上非零常数 \(2^{-n}\) 所得，因而可逆；切片张成整个 \(N+1\) 维多项式空间。若它有 \(r\) 项分离表示，全部切片都在至多 \(r\) 维的 \(\operatorname{span}\{f_i\}\) 中，故 \(r\ge N+1\)。题述展开恰有 \(N+1\) 项，达到该下界。
\end{solution}
